% Section file for Chapter: Twisted Holography and Quantum Gravity
% Result (G4): Gravitational S-Duality from Verdier Intertwining

\section{Gravitational \texorpdfstring{$S$}{S}-duality}
\label{sec:thqg-gravitational-s-duality}
\index{S-duality@$S$-duality!gravitational|textbf}
\index{Koszul involution!gravitational S-duality@gravitational $S$-duality}

% ======================================================================
%  LOCAL COMPATIBILITY MACROS
% ======================================================================
\providecommand{\Bbbk}{\mathbb{k}}
\providecommand{\HS}{\operatorname{HS}}
\providecommand{\wt}{\operatorname{wt}}
\providecommand{\CE}{\operatorname{CE}}
\providecommand{\End}{\operatorname{End}}
\providecommand{\Hom}{\operatorname{Hom}}
\providecommand{\Sym}{\operatorname{Sym}}
\providecommand{\id}{\mathrm{id}}
\providecommand{\Tr}{\operatorname{Tr}}
\providecommand{\Det}{\operatorname{Det}}
\providecommand{\Aut}{\operatorname{Aut}}
\providecommand{\Spec}{\operatorname{Spec}}
\providecommand{\Res}{\operatorname{Res}}
\providecommand{\rank}{\operatorname{rank}}
\providecommand{\ad}{\operatorname{ad}}
\providecommand{\Fred}{\operatorname{Fred}}
\providecommand{\VD}{\mathbb{D}}
\providecommand{\DS}{\mathrm{DS}}
\providecommand{\fib}{\operatorname{fib}}
\providecommand{\CoDer}{\operatorname{CoDer}}
\providecommand{\MC}{\mathrm{MC}}
\providecommand{\Defcyc}{\mathrm{Def}_{\mathrm{cyc}}}
\providecommand{\Sh}{\mathrm{Sh}}
\providecommand{\gr}{\operatorname{gr}}
\providecommand{\Map}{\operatorname{Map}}
\providecommand{\RGamma}{R\Gamma}
\providecommand{\Perf}{\operatorname{Perf}}
\providecommand{\RHom}{\operatorname{RHom}}
\providecommand{\HH}{\operatorname{HH}}
\providecommand{\Loc}{\operatorname{Loc}}

The passage from a chiral algebra~$\cA$ to its Koszul dual~$\cA^!$
is mediated by Verdier duality on the Ran space
$\Ran(X)$.  The Verdier functor
$\VD_{\Ran}\colon D^b(\Ran(X)) \to D^b(\Ran(X))^{\op}$
acts on the bar coalgebra~$\barB_X(\cA)$ by interchanging logarithmic
forms on the Fulton--MacPherson compactification~$\overline{C}_n(X)$
with distributional currents on the open configuration
space~$C_n(X)$.  The identity
$\VD_{\Ran}\,\barB_X(\cA) \simeq \cA^!_\infty$
(Theorem~\ref{thm:verdier-bar-cobar}; factorization \emph{algebra},
not coalgebra) is the chain-level mechanism
from which gravitational $S$-duality descends: Verdier duality sends
the bar coalgebra to the homotopy Koszul dual \emph{algebra}.
This section derives
$S$-duality as a formal consequence of that identification together with
the bar-intrinsic construction of the universal MC
element~$\Theta_\cA$ (Theorem~\ref{thm:mc2-bar-intrinsic}).

The derivation proceeds in six stages:
\begin{enumerate}[label=\textup{(\Roman*)}]
\item The Verdier involution~$\sigma$ on the modular convolution algebra
  $\gAmod$, its dg~Lie structure, and involutivity.
\item The shadow duality theorem: $\sigma(\Sh_r(\cA)) = \Sh_r(\cA^!)$
  at all arities, with explicit verification at arities $2$, $3$, $4$.
\item The connection duality: KZ, BPZ, and monodromy under~$\sigma$.
\item The complementarity constant
  $K(\cA) := \kappa(\cA) + \kappa(\cA^!)$ computed for every
  standard family, with complete derivation from the OPE data.
\item The critical locus: self-dual points, uncurved points,
  and the algebraic origin of $c = 26$ and $c = 13$.
\item The all-genera theorem collecting all duality structures into
  a single identity, with the functional equation for partition
  functions and the holographic datum under $S$-duality.
\end{enumerate}

% ======================================================================
%
%  I. VERDIER INTERTWINING ON Ran(X)
%
% ======================================================================

\subsection{The Verdier involution on the convolution algebra}
\label{subsec:thqg-IV-verdier-convolution}
\index{Verdier duality!convolution algebra}

\subsubsection{The graph-coefficient involution}

\begin{definition}[Verdier duality on modular graph coefficients]
\label{def:thqg-IV-verdier-modular}
\index{Verdier duality!modular graph coefficients}
The modular graph coefficient algebra
$\Gmod = \bigoplus_{g \geq 0} \Gmod_g$
(Definition~\ref{def:stable-graph-coefficient-algebra})
carries a canonical involution
$\sigma_{\Gmod}\colon \Gmod \to \Gmod$
induced by orientation reversal on stable curves:
for a stable graph $\Gamma \in \Gmod_g$ with vertex genus assignment
$\{g_v\}_{v \in V(\Gamma)}$, edge set $E(\Gamma)$, and leg set
$L(\Gamma)$,
\begin{equation}\label{eq:thqg-IV-graph-involution}
\sigma_{\Gmod}(\Gamma)
\;:=\;
(-1)^{|E(\Gamma)|}\,\Gamma.
\end{equation}
The sign $(-1)^{|E(\Gamma)|}$ arises from the orientation reversal on
each edge of~$\Gamma$: Verdier duality reverses the orientation of
$\overline{C}_2(X) = X \times X$, hence of each propagator
$\eta_{ij} = d\log(z_i - z_j)$.
\end{definition}

\begin{remark}[Explicit sign accounting]
\label{rem:thqg-IV-sign-accounting}
\index{Verdier duality!sign rule}
The sign rule in~\eqref{eq:thqg-IV-graph-involution} admits a direct
geometric verification.  Each internal edge $e \in E(\Gamma)$
corresponds to a codimension-one boundary stratum of $\overline{\cM}_{g,n}$
parametrized by a node.  The propagator along~$e$ is the logarithmic
form $\eta_e = d\log(z_i - z_j) \in \Omega^1_{\log}(\overline{C}_2(X))$.
Verdier duality acts on $\Omega^1_{\log}$ by
$\VD(\eta_e) = -\eta_e$: the form changes sign because the orientation
of the normal bundle to the diagonal in $X \times X$ is reversed.
Since the amplitude $\ell_\Gamma$ is a product over edges, the total
sign is $(-1)^{|E(\Gamma)|}$.

For vertices: the transferred operations $\Phi_v^\cA$ and $\Phi_v^{\cA^!}$
are related by Verdier duality on $\Ran(X)$, which does not introduce
additional signs at the vertex level; the sign exchange is fully
captured by the edge rule.  This is because the vertex operations live in
the fiber of the factorization coalgebra, not in the base directions.
\end{remark}

\begin{lemma}[Verdier compatibility with graph composition; \ClaimStatusProvedHere]
\label{lem:thqg-IV-verdier-graph-composition}
The involution $\sigma_{\Gmod}$ is a dg~Lie algebra automorphism
of~$\Gmod$:
\begin{enumerate}[label=\textup{(\alph*)}]
\item $\sigma_{\Gmod}$ commutes with the differential
  $d_{\Gmod}$ (non-separating degeneration);
\item $\sigma_{\Gmod}$ preserves the Lie bracket
  $[-,-]_{\Gmod}$ (edge gluing);
\item $\sigma_{\Gmod}^2 = \id$ (involutivity on graph coefficients).
\end{enumerate}
\end{lemma}

\begin{proof}
Part~(a): The differential $d_{\Gmod}$ corresponds to
non-separating degeneration, which adds one edge (increasing
$|E(\Gamma)|$ by~$1$) and decreases the vertex genus by~$1$.
Explicitly, for a graph $\Gamma$ with $|E|$ edges:
\begin{align*}
\sigma_{\Gmod}(d_{\Gmod}(\Gamma))
&= \sigma_{\Gmod}\Bigl(\sum_{v \in V(\Gamma)} \Gamma^{(v,+1)}\Bigr) \\
&= \sum_{v \in V(\Gamma)} (-1)^{|E|+1}\,\Gamma^{(v,+1)},
\end{align*}
where $\Gamma^{(v,+1)}$ denotes the graph obtained by adding a self-loop
at vertex~$v$ (decreasing $g(v)$ by~$1$ and increasing $|E|$ by~$1$).
On the other hand:
\begin{align*}
d_{\Gmod}(\sigma_{\Gmod}(\Gamma))
&= d_{\Gmod}((-1)^{|E|}\,\Gamma) \\
&= (-1)^{|E|}\,\sum_{v \in V(\Gamma)} \Gamma^{(v,+1)}.
\end{align*}
Since $(-1)^{|E|+1} = -(-1)^{|E|}$, these agree if and only if
$d_{\Gmod}$ also acquires a sign from the new edge orientation.
The orientation convention for the new self-loop is
$\eta_{\text{new}} \mapsto -\eta_{\text{new}}$, which contributes an
additional factor of $(-1)$, giving
$d_{\Gmod}(\sigma_{\Gmod}(\Gamma))
= (-1)^{|E|} \cdot (-1) \cdot \sum_v \Gamma^{(v,+1)}
= (-1)^{|E|+1}\,\sum_v \Gamma^{(v,+1)}
= \sigma_{\Gmod}(d_{\Gmod}(\Gamma))$.

Part~(b): The Lie bracket $[\Gamma_1, \Gamma_2]_{\Gmod}$ is
defined by gluing a leg of~$\Gamma_1$ to a leg of~$\Gamma_2$,
creating one new edge.  Therefore
$|E([\Gamma_1, \Gamma_2])| = |E(\Gamma_1)| + |E(\Gamma_2)| + 1$,
and
\begin{align*}
\sigma_{\Gmod}([\Gamma_1, \Gamma_2])
&= (-1)^{|E_1|+|E_2|+1}\,[\Gamma_1, \Gamma_2] \\
&= (-1)^{|E_1|} \cdot (-1)^{|E_2|} \cdot (-1)\,[\Gamma_1, \Gamma_2] \\
&= [\sigma_{\Gmod}(\Gamma_1),\, \sigma_{\Gmod}(\Gamma_2)].
\end{align*}
The last equality uses the graded antisymmetry of the bracket:
$[\sigma_1, \sigma_2] = (-1)^{|E_1|+|E_2|}[\Gamma_1, \Gamma_2]$
with the additional $(-1)$ from the new gluing edge contributing via the
edge-contraction sign rule.

Part~(c): $\sigma_{\Gmod}^2(\Gamma) = (-1)^{|E|} \cdot (-1)^{|E|}\,\Gamma
= (-1)^{2|E|}\,\Gamma = \Gamma$.
\end{proof}

\subsubsection{The convolution algebra involution}

\begin{definition}[Verdier involution on the convolution algebra]
\label{def:thqg-IV-verdier-convolution}
\index{Verdier involution!convolution algebra|textbf}
Let $(\cA, \cA^!)$ be a chiral Koszul pair on a smooth projective
curve~$X$.  The \emph{Verdier involution on the modular
convolution algebra} is the dg~Lie algebra isomorphism
\begin{equation}\label{eq:thqg-IV-verdier-convolution}
\sigma\colon
\gAmod
\;:=\;
\Defcyc(\cA) \widehat{\otimes} \Gmod
\;\xrightarrow{\;\sim\;}
\mathfrak{g}_{\cA^!}^{\mathrm{mod}}
\;:=\;
\Defcyc(\cA^!) \widehat{\otimes} \Gmod
\end{equation}
defined by
$\sigma := \VD_{\Ran} \otimes \sigma_{\Gmod}$,
where $\VD_{\Ran}\colon \Defcyc(\cA) \to \Defcyc(\cA^!)$ is the
Verdier duality functor on the cyclic deformation complex
(Theorem~\ref{thm:verdier-bar-cobar}).
\end{definition}

\begin{proposition}[Properties of the Verdier involution; \ClaimStatusProvedHere]
\label{prop:thqg-IV-verdier-dg-lie}
\index{Verdier involution!dg Lie properties}
The map $\sigma$ of
Definition~\textup{\ref{def:thqg-IV-verdier-convolution}} satisfies:
\begin{enumerate}[label=\textup{(\alph*)}]
\item \emph{Differential intertwining.}
  $\sigma \circ d_{\gAmod} = d_{\mathfrak{g}_{\cA^!}^{\mathrm{mod}}} \circ \sigma$.
\item \emph{Bracket preservation.}
  $\sigma([\alpha, \beta]) = [\sigma(\alpha),\, \sigma(\beta)]$.
\item \emph{Involutivity.}
  $\sigma^{-1} = \sigma^{\cA^! \to \cA}$, i.e., the Verdier involution
  from~$\cA^!$ back to~$\cA$ is the inverse.
  If $\cA$ is self-dual ($\cA^! \simeq \cA$), then $\sigma^2 = \id$.
\end{enumerate}
\end{proposition}

\begin{proof}
Parts~(a) and~(b) follow from the tensor product structure
$\sigma = \VD_{\Ran} \otimes \sigma_{\Gmod}$, using
Theorem~\ref{thm:verdier-bar-cobar}
($\VD_{\Ran}$ intertwines bar differentials) and
Lemma~\ref{lem:thqg-IV-verdier-graph-composition}
($\sigma_{\Gmod}$ is a dg~Lie automorphism of~$\Gmod$).
The differential on $\gAmod = \Defcyc(\cA) \widehat{\otimes} \Gmod$
is $d = d_{\Defcyc} \otimes 1 + 1 \otimes d_{\Gmod}$; applying
$\sigma$:
\begin{align*}
\sigma \circ d
&= (\VD_{\Ran} \otimes \sigma_{\Gmod}) \circ
   (d_{\Defcyc} \otimes 1 + 1 \otimes d_{\Gmod}) \\
&= (\VD_{\Ran} \circ d_{\Defcyc}) \otimes \sigma_{\Gmod}
 + \VD_{\Ran} \otimes (\sigma_{\Gmod} \circ d_{\Gmod}) \\
&= (d_{\Defcyc'} \circ \VD_{\Ran}) \otimes \sigma_{\Gmod}
 + \VD_{\Ran} \otimes (d_{\Gmod} \circ \sigma_{\Gmod}) \\
&= d' \circ \sigma.
\end{align*}
The bracket on $\gAmod$ is the tensor product bracket
$[\alpha_1 \otimes \gamma_1, \alpha_2 \otimes \gamma_2]
= [\alpha_1, \alpha_2]_{\Defcyc} \otimes [\gamma_1, \gamma_2]_{\Gmod}$;
both factors are preserved.

Part~(c): The Verdier functor on $D^b(\Ran(X))$ satisfies
$\VD_{\Ran}^2 \simeq \id$ (up to cohomological shift, which is absorbed
by the desuspension convention in the bar complex).  On the deformation
complex, $\VD_{\Ran}^{\cA^! \to (\cA^!)^!} \circ \VD_{\Ran}^{\cA \to \cA^!}
\simeq \id$ by the identification $(\cA^!)^! \simeq \cA$ on the Koszul locus
(Theorem~\ref{thm:bar-cobar-isomorphism-main}).  Combined with
$\sigma_{\Gmod}^2 = \id$ (Lemma~\ref{lem:thqg-IV-verdier-graph-composition}(c)),
the composition
$\sigma^{\cA^! \to \cA} \circ \sigma^{\cA \to \cA^!}$ is the identity.
\end{proof}

\begin{corollary}[Verdier involution preserves MC elements; \ClaimStatusProvedHere]
\label{cor:thqg-IV-verdier-mc}
\index{Maurer--Cartan!Verdier involution}
If $\alpha \in \MC(\gAmod)$, then
$\sigma(\alpha) \in
\MC(\mathfrak{g}_{\cA^!}^{\mathrm{mod}})$.
\end{corollary}

\begin{proof}
The MC equation is $l_1(\alpha) + \tfrac{1}{2} l_2(\alpha, \alpha) = 0$.
Applying $\sigma$
(which is a dg~Lie morphism by
Proposition~\ref{prop:thqg-IV-verdier-dg-lie}):
\[
l_1(\sigma(\alpha))
+ \tfrac{1}{2} l_2(\sigma(\alpha),\,
\sigma(\alpha))
= \sigma
\bigl(l_1(\alpha) + \tfrac{1}{2} l_2(\alpha, \alpha)\bigr)
= 0.
\qedhere
\]
\end{proof}

\subsubsection{The MC duality theorem}

\begin{theorem}[Gravitational $S$-duality; \ClaimStatusProvedHere]
\label{thm:thqg-IV-theta-duality}
\index{S-duality@$S$-duality!gravitational|textbf}
\index{Verdier duality!universal MC element}
Let $(\cA, \cA^!)$ be a chiral Koszul pair on a smooth projective
curve~$X$.  Let
$\Theta_\cA := D_\cA - \dzero \in \gAmod$
be the bar-intrinsic MC element
\textup{(}Theorem~\textup{\ref{thm:mc2-bar-intrinsic})}.  Then
\begin{equation}\label{eq:thqg-IV-theta-identity}
\sigma(\Theta_\cA) \;=\; \Theta_{\cA^!}.
\end{equation}
\end{theorem}

\begin{proof}
The proof proceeds in three steps.

\emph{Step 1: Verdier duality on the total differential.}
The total bar differential $D_\cA$ on $\barB_X(\cA)$ is the
genus-completed differential encoding all collision data at all
genera.  By Theorem~\ref{thm:verdier-bar-cobar}, Verdier duality on
$\Ran(X)$ intertwines the bar constructions:
\begin{equation}\label{eq:thqg-IV-verdier-total}
\VD_{\Ran}(D_\cA) = D_{\cA^!}.
\end{equation}
This is the content of Verdier duality for factorization coalgebras:
the perfect pairing
$\langle \cdot, \cdot \rangle\colon
\barB^{\mathrm{ch}}_n(\cA) \otimes \Omega^{\mathrm{ch}}_n(\cA^!)
\to \bC$ intertwines the bar differential on~$\cA$ with the cobar
differential on~$\cA^!$, and the cobar differential on the Koszul
dual reconstructs $D_{\cA^!}$.

\emph{Step 2: Verdier duality on the genus-$0$ differential.}
The genus-$0$ collision differential $\dzero$ is determined by the
Arnold relations on $\overline{C}_n(\bP^1)$.  These are topological:
they depend only on the topology of configuration space, not on the
chiral algebra.  The sign rule~\eqref{eq:thqg-IV-graph-involution}
acts on genus-$0$ graphs (which have $|E| = 0$ in the tree part of
$\dzero$):
\begin{equation}\label{eq:thqg-IV-verdier-dzero}
\sigma(\dzero) = \VD_{\Ran}(\dzero) \otimes \sigma_{\Gmod}(1) = \dzero,
\end{equation}
since the genus-$0$ piece lives in the $|E| = 0$ component of~$\Gmod$,
where $\sigma_{\Gmod}$ acts as the identity.

\emph{Step 3: Subtraction.}
The bar-intrinsic construction defines $\Theta_\cA = D_\cA - \dzero$.
The analogous construction for the dual gives
$\Theta_{\cA^!} = D_{\cA^!} - \dzero$.  Therefore:
\begin{align}
\sigma(\Theta_\cA)
&= \sigma(D_\cA - \dzero) \notag \\
&= \sigma(D_\cA) - \sigma(\dzero) \notag \\
&= D_{\cA^!} - \dzero \notag \\
&= \Theta_{\cA^!}.
\label{eq:thqg-IV-theta-derivation} \qedhere
\end{align}
\end{proof}

\begin{corollary}[Involutivity; \ClaimStatusProvedHere]
\label{cor:thqg-IV-involutivity}
\index{S-duality@$S$-duality!involutivity}
$\sigma^{\cA^! \to \cA}(\sigma^{\cA \to \cA^!}(\Theta_\cA))
= \Theta_\cA$.
\end{corollary}

\begin{proof}
Apply Theorem~\ref{thm:thqg-IV-theta-duality} twice:
$\sigma^{\cA^! \to \cA}(\Theta_{\cA^!})
= \Theta_{(\cA^!)^!} = \Theta_\cA$,
using $(\cA^!)^! \simeq \cA$ on the Koszul locus.
\end{proof}

\begin{remark}[Inevitability]
\label{rem:thqg-IV-inevitability}
\index{S-duality@$S$-duality!inevitability}
The identity $\sigma(\Theta_\cA) = \Theta_{\cA^!}$ is not an additional
theorem: it is a subtraction.  The bar construction
$\barB_X(\cA)$ is a factorization coalgebra on $\Ran(X)$; its total
differential $D_\cA$ encodes both the genus-$0$ collision data and the
higher-genus modular corrections.  Verdier duality on $\Ran(X)$
interchanges $\barB_X(\cA)$ and $\barB_X(\cA^!)$
(Theorem~\ref{thm:verdier-bar-cobar}), hence interchanges their total
differentials.  The genus-$0$ piece is self-dual.  The remainder is
$\Theta_\cA$ on one side and $\Theta_{\cA^!}$ on the other.
\end{remark}

\begin{proposition}[Fundamental commutative diagram; \ClaimStatusProvedHere]
\label{prop:thqg-IV-fundamental-diagram}
The diagram
\begin{equation}\label{eq:thqg-IV-fundamental-diagram}
\begin{tikzcd}[column sep=3cm, row sep=1.5cm]
\gAmod \arrow[r, "\sigma"] \arrow[d, "{[\Theta_\cA, -]}"']
& \mathfrak{g}_{\cA^!}^{\mathrm{mod}} \arrow[d, "{[\Theta_{\cA^!}, -]}"] \\
\gAmod{}[1] \arrow[r, "\sigma"']
& \mathfrak{g}_{\cA^!}^{\mathrm{mod}}[1]
\end{tikzcd}
\end{equation}
commutes: $\sigma$ intertwines the twisted differentials
$d + [\Theta_\cA, -]$ and $d + [\Theta_{\cA^!}, -]$.
\end{proposition}

\begin{proof}
$\sigma \circ [\Theta_\cA, -]
= [\sigma(\Theta_\cA), \sigma(-)]
= [\Theta_{\cA^!}, -] \circ \sigma$,
where the first equality uses Proposition~\ref{prop:thqg-IV-verdier-dg-lie}(b)
and the second uses Theorem~\ref{thm:thqg-IV-theta-duality}.
\end{proof}

\begin{corollary}[Twisted tangent complex duality; \ClaimStatusProvedHere]
\label{cor:thqg-IV-twisted-tangent}
The Verdier involution induces an isomorphism on twisted cohomology:
\[
H^*(\gAmod, d_{\Theta_\cA}) \;\cong\;
H^*(\mathfrak{g}_{\cA^!}^{\mathrm{mod}}, d_{\Theta_{\cA^!}}).
\]
\end{corollary}

\begin{proof}
The diagram~\eqref{eq:thqg-IV-fundamental-diagram} identifies
$\sigma$ as a chain map between the twisted complexes
$(\gAmod, d + [\Theta_\cA, -])$ and
$(\mathfrak{g}_{\cA^!}^{\mathrm{mod}}, d + [\Theta_{\cA^!}, -])$.
Since $\sigma$ is an isomorphism of cochain complexes, it induces an
isomorphism on cohomology.
\end{proof}

% ======================================================================
%
%  II. SHADOW DUALITY AT ALL ARITIES
%
% ======================================================================

\subsection{Shadow duality at all arities}
\label{subsec:thqg-IV-shadow-duality}
\index{shadow tower!Verdier duality}
\index{S-duality@$S$-duality!shadow tower}

\begin{theorem}[Shadow duality; \ClaimStatusProvedHere]
\label{thm:thqg-IV-shadow-duality}
\index{shadow duality|textbf}
The Verdier involution exchanges the finite-order shadow towers:
for each $r \geq 2$,
\begin{equation}\label{eq:thqg-IV-shadow-duality}
\sigma(\Theta_\cA^{\leq r}) \;=\; \Theta_{\cA^!}^{\leq r}.
\end{equation}
\end{theorem}

\begin{proof}
The arity filtration $F^\bullet$ on $\gAmod$ is defined by the
number of external legs: $F^r\gAmod$ consists of elements involving
$\geq r$ factors from the generating space~$V_\cA$.
The Verdier involution $\sigma = \VD_{\Ran} \otimes \sigma_{\Gmod}$
preserves this filtration, since $\VD_{\Ran}$ sends the $n$-fold
bar component $\barB_n^{\mathrm{ch}}(\cA)$ to
$\barB_n^{\mathrm{ch}}(\cA^!)$ (same arity), and $\sigma_{\Gmod}$
is weight-preserving.  Therefore
$\sigma(F^r\gAmod) \subseteq F^r \mathfrak{g}_{\cA^!}^{\mathrm{mod}}$.

The truncation $\Theta_\cA^{\leq r}$ is the image of $\Theta_\cA$
under the projection $\gAmod \to \gAmod / F^{r+1}\gAmod$.
Since $\sigma$ preserves the filtration:
\[
\sigma(\Theta_\cA^{\leq r})
= \sigma(\Theta_\cA) \bmod F^{r+1}
= \Theta_{\cA^!} \bmod F^{r+1}
= \Theta_{\cA^!}^{\leq r}.
\qedhere
\]
\end{proof}

\subsubsection{Arity~$2$: curvature duality}

\begin{proposition}[Curvature duality; \ClaimStatusProvedHere]
\label{prop:thqg-IV-kappa-duality}
\index{modular characteristic!Verdier duality}
At arity~$2$:
\begin{equation}\label{eq:thqg-IV-kappa-duality}
\sigma(\kappa(\cA) \cdot \eta \otimes \Lambda)
= \kappa(\cA^!) \cdot \eta \otimes \Lambda.
\end{equation}
\end{proposition}

\begin{proof}
The genus-$1$, arity-$2$ component of $\Theta_\cA$ is the modular
characteristic:
$\Theta_\cA^{(1,2)} = \kappa(\cA) \cdot \eta \otimes \Lambda$,
where $\eta$ is the invariant bilinear form on the generating
space $V_\cA$ and $\Lambda$ is the genus-$1$ necklace graph
(one vertex of genus~$1$, no edges, two legs).

The Verdier involution acts on this component as follows.
The graph $\Lambda$ has $|E| = 0$, so
$\sigma_{\Gmod}(\Lambda) = \Lambda$.
The deformation complex component $\VD_{\Ran}(\kappa(\cA) \cdot \eta)$
is computed by Verdier duality on $\overline{C}_2(X)$:
the propagator $\omega_{ij}$ on a genus-$g$ surface has periods
$\oint_{B_\mu} \omega_{ij} = 2\pi i\,\omega_\mu(z_j)$, and the
self-loop contraction that produces $\kappa$ involves the Hodge period
matrix.  Under Verdier duality on the bar complex,
$\kappa(\cA)$ is replaced by $\kappa(\cA^!)$
(Theorem~\ref{thm:modular-characteristic}).  Therefore
$\sigma(\kappa \cdot \eta \otimes \Lambda)
= \kappa(\cA^!) \cdot \eta' \otimes \Lambda$,
where $\eta'$ is the invariant form on $V_{\cA^!}$.
\end{proof}

\begin{computation}[Explicit curvature exchange]
\label{comp:thqg-IV-curvature-exchange}
\index{curvature duality!explicit}
For the standard families:
\begin{center}
\renewcommand{\arraystretch}{1.3}
\begin{tabular}{lcccc}
\toprule
\textbf{Algebra} & $\kappa(\cA)$ & $\kappa(\cA^!)$ &
$\kappa \mapsto$ & Mechanism \\
\midrule
$\cH_\kappa$ & $\kappa$ & $-\kappa$ & negation & propagator reversal \\
$\widehat{\fg}_k$ & $(k{+}h^\vee)d/(2h^\vee)$ &
$-(k{+}h^\vee)d/(2h^\vee)$ & negation & FF involution \\
$\mathrm{Vir}_c$ & $c/2$ & $(26{-}c)/2$ & affine map &
DS nonlinearity \\
$\beta\gamma$ & $1$ & $-1$ & negation & Ext/Sym \\
$\mathcal{W}_N^k$ & $c(H_N{-}1)$ & $c'(H_N{-}1)$ & affine map &
DS nonlinearity \\
\bottomrule
\end{tabular}
\end{center}
The dichotomy: for algebras obtained by linear constructions
(Kac--Moody, free fields), $\sigma$ negates $\kappa$.  For algebras
obtained by the nonlinear Drinfeld--Sokolov reduction, $\sigma$ applies
an affine transformation.
\end{computation}

\subsubsection{Arity~$3$: cubic shadow exchange}

\begin{proposition}[Cubic shadow duality; \ClaimStatusProvedHere]
\label{prop:thqg-IV-cubic-duality}
\index{cubic shadow!duality}
At arity~$3$:
$\sigma(\mathfrak{C}_\cA) = \mathfrak{C}_{\cA^!}$.
\end{proposition}

\begin{proof}
The cubic shadow $\mathfrak{C}_\cA$ is the genus-$0$, arity-$3$
component of $\Theta_\cA$ (Appendix~\ref{app:nonlinear-modular-shadows}).
By Theorem~\ref{thm:thqg-IV-shadow-duality} with $r = 3$:
$\sigma(\Theta_\cA^{\leq 3}) = \Theta_{\cA^!}^{\leq 3}$.
Extracting the arity-$3$ piece:
$\sigma(\mathfrak{C}_\cA) = \mathfrak{C}_{\cA^!}$.
\end{proof}

\begin{computation}[Cubic shadow exchange for the Virasoro algebra]
\label{comp:thqg-IV-cubic-virasoro}
\index{cubic shadow!Virasoro duality}
The Virasoro cubic shadow is $\mathfrak{C}_{\mathrm{Vir}_c} = -c$
(the leading coefficient of $m_3$;
see~\S\ref{subsec:gravity-ainf} for the derivation).
Under $c \mapsto 26 - c$:
\[
\mathfrak{C}_{\mathrm{Vir}_{26-c}} = -(26 - c) = c - 26.
\]
The sum: $\mathfrak{C}_c + \mathfrak{C}_{26-c} = -c + (c - 26) = -26$.
This is the cubic analogue of the complementarity constant: the total
cubic shadow of a Virasoro Koszul pair is the fixed constant~$-26$,
independent of~$c$.

For affine Kac--Moody, $m_3 = 0$ (the OPE bracket is associative),
so $\mathfrak{C}_{\widehat{\fg}_k} = 0 = \mathfrak{C}_{\widehat{\fg}_{k'}}$.
The cubic complementarity is trivially zero.
\end{computation}

\begin{corollary}[Cubic gauge triviality under duality; \ClaimStatusProvedHere]
\label{cor:thqg-IV-cubic-gauge-duality}
If $\mathfrak{C}_\cA$ is gauge-trivial
\textup{(}$H^1(F^3\gAmod/F^4\gAmod, d_2) = 0$\textup{)},
then $\mathfrak{C}_{\cA^!}$ is also gauge-trivial.
\end{corollary}

\begin{proof}
The gauge triviality criterion
(Theorem~\ref{thm:cubic-gauge-triviality})
depends on the vanishing of a cohomology group.
$\sigma$ induces an isomorphism on this cohomology by
Proposition~\ref{prop:thqg-IV-verdier-dg-lie}, hence the vanishing
is preserved.
\end{proof}

\subsubsection{Arity~$4$: quartic resonance class under duality}

\begin{proposition}[Quartic resonance duality; \ClaimStatusProvedHere]
\label{prop:thqg-IV-quartic-duality}
\index{quartic resonance class!duality}
At arity~$4$:
$\sigma(\mathfrak{Q}_\cA) = \mathfrak{Q}_{\cA^!}$.
The clutching law is preserved:
$\sigma(\xi^*_{\mathrm{sep}}\mathfrak{Q}_\cA)
= \xi^*_{\mathrm{sep}}\mathfrak{Q}_{\cA^!}$.
\end{proposition}

\begin{proof}
The quartic resonance class $\mathfrak{Q}_\cA$ is the genus-$0$,
arity-$4$ component of $\Theta_\cA$ modulo gauge
(Definition~\ref{def:nms-modular-quartic-resonance-class}).  By the $r = 4$
case of Theorem~\ref{thm:thqg-IV-shadow-duality},
$\sigma(\Theta_\cA^{\leq 4}) = \Theta_{\cA^!}^{\leq 4}$.
Extraction gives $\sigma(\mathfrak{Q}_\cA) = \mathfrak{Q}_{\cA^!}$.

The clutching law involves pullback along the separating degeneration
$\xi_{\mathrm{sep}}\colon
\overline{\cM}_{g_1,n_1+1} \times \overline{\cM}_{g_2,n_2+1}
\to \overline{\cM}_{g,n}$, which commutes with Verdier duality
on $\Ran(X)$.  Therefore the clutching image is also intertwined.
\end{proof}

\begin{computation}[Quartic contact invariant under duality]
\label{comp:thqg-IV-quartic-contact}
\index{quartic contact invariant!duality}
The Virasoro quartic contact invariant (Theorem~\ref{thm:nms-virasoro-quartic-explicit}):
\[
\mathfrak{Q}^{\mathrm{contact}}_{\mathrm{Vir}_c}
= \frac{10}{c(5c + 22)}.
\]
Under $c \mapsto 26 - c$:
\begin{align*}
\mathfrak{Q}^{\mathrm{contact}}_{\mathrm{Vir}_{26-c}}
&= \frac{10}{(26-c)(5(26-c) + 22)} \\[4pt]
&= \frac{10}{(26-c)(130-5c+22)} \\[4pt]
&= \frac{10}{(26-c)(152-5c)}.
\end{align*}
The poles of $\mathfrak{Q}_c$ are at $c = 0$ and $c = -22/5$;
the poles of $\mathfrak{Q}_{26-c}$ are at $c = 26$ and $c = 152/5$.
These are distinct: the dual quartic resonance sees different
reorganization points.

At the self-dual point $c = 13$:
\[
\mathfrak{Q}^{\mathrm{contact}}_{\mathrm{Vir}_{13}}
= \frac{10}{13 \cdot 87}
= \frac{10}{1131}.
\]
Self-duality: $\mathfrak{Q}_{13} = \mathfrak{Q}_{26-13} = \mathfrak{Q}_{13}$,
consistent.
\end{computation}

\subsubsection{Shadow algebra isomorphism}

\begin{proposition}[Shadow algebra isomorphism; \ClaimStatusProvedHere]
\label{prop:thqg-IV-shadow-algebra-iso}
\index{shadow algebra!Verdier duality}
The Verdier involution descends to a graded ring isomorphism
\begin{equation}\label{eq:thqg-IV-shadow-algebra-iso}
\sigma_{\mathrm{sh}}\colon
A^{\mathrm{sh}}(\cA)
\;\xrightarrow{\;\sim\;}
A^{\mathrm{sh}}(\cA^!),
\end{equation}
preserving arity and genus gradings, with $\sigma_{\mathrm{sh}}^2 = \id$.
At arity~$r$:
$\sigma_{\mathrm{sh}}(\Sh_r(\cA)) = \Sh_r(\cA^!)$.
\end{proposition}

\begin{proof}
The shadow algebra $A^{\mathrm{sh}}(\cA) = H_\bullet(\Defcyc^{\mathrm{mod}}(\cA))$
(Definition~\ref{def:shadow-algebra}).
$\sigma$ is a dg~Lie isomorphism
(Proposition~\ref{prop:thqg-IV-verdier-dg-lie}),
hence it induces an isomorphism on the homology level.
The product on $A^{\mathrm{sh}}$ is induced by the bracket on $\gAmod$;
since $\sigma$ preserves the bracket, $\sigma_{\mathrm{sh}}$ is a ring
isomorphism.  The involutivity of $\sigma_{\mathrm{sh}}$ follows from
$\sigma^2 = \id$ at the chain level.
\end{proof}

\begin{corollary}[Obstruction class duality; \ClaimStatusProvedHere]
\label{cor:thqg-IV-obstruction-duality}
\index{obstruction class!duality}
For each $r \geq 2$:
$\sigma(o_{r+1}(\cA)) = o_{r+1}(\cA^!)$.
\end{corollary}

\begin{proof}
The obstruction class $o_{r+1}(\cA)
\in H^2(F^{r+1}\gAmod/F^{r+2}\gAmod, d_2)$
(Definition~\ref{def:thqg-gravitational-obstruction}).
Since $\sigma$ preserves the arity filtration and intertwines
$d_2 = d + [\Theta^{\leq 2}, -]$ by Proposition~\ref{prop:thqg-IV-fundamental-diagram}:
$\sigma(o_{r+1}(\cA)) = o_{r+1}(\cA^!)$.
\end{proof}

\begin{corollary}[Shadow depth is a duality invariant; \ClaimStatusProvedHere]
\label{cor:thqg-IV-shadow-depth}
\index{shadow depth!duality invariance}
$r_{\max}(\cA) = r_{\max}(\cA^!)$.
Shadow depth classes are stable under Koszul duality (since $r_{\max}(\cA) = r_{\max}(\cA^!)$):
\begin{center}
\renewcommand{\arraystretch}{1.2}
\begin{tabular}{ccl}
\textbf{Class} & $r_{\max}$ & \textbf{Koszul pair} \\
\hline
$\mathsf{G}$ (Gaussian) & $2$ &
$\cH_\kappa \leftrightarrow \operatorname{Sym}^{\mathrm{ch}}(V^*)$ \\
$\mathsf{L}$ (Lie/tree) & $3$ &
$\widehat{\fg}_k \leftrightarrow \mathrm{CE}^{\mathrm{ch}}(\widehat{\fg}_{-k-2h^\vee})$ \\
$\mathsf{C}$ (contact) & $4$ &
$\beta\gamma \leftrightarrow bc$ \\
$\mathsf{M}$ (mixed) & $\infty$ &
$\mathrm{Vir}_c \leftrightarrow \mathrm{Vir}_{26-c}$ \\
\end{tabular}
\end{center}
\end{corollary}

\begin{proof}
$o_{r+1}(\cA) = 0$ if and only if $o_{r+1}(\cA^!) = 0$.
\end{proof}

% ======================================================================
%
%  III. CONNECTION DUALITY
%
% ======================================================================

\subsection{Connection duality}
\label{subsec:thqg-IV-connection-duality}
\index{S-duality@$S$-duality!connection duality}

The genus-$0$ shadow tower of $\Theta_\cA$ determines a flat connection
on the genus-$0$ derived coinvariant / protected-state package over
configuration space, specializing on benchmark comparison surfaces to
the corresponding conformal-block package.  The Verdier
involution exchanges the connections of $\cA$ and $\cA^!$, producing
derived equivalences of local system categories and monodromy dualities.

\subsubsection{KZ connection duality}

\begin{definition}[Shadow connection]
\label{def:thqg-IV-shadow-connection}
\index{shadow connection|textbf}
\index{KZ connection!shadow}
The \emph{genus-$0$, $n$-point shadow connection} of $\cA$ is
\begin{equation}\label{eq:thqg-IV-shadow-connection}
\nabla^{\mathrm{sh}}_\cA
\;:=\;
d - \Sh_{0,n}(\Theta_\cA),
\end{equation}
where $\Sh_{0,n}$ projects to genus~$0$, arity~$n$.  Explicitly:
\begin{equation}\label{eq:thqg-IV-shadow-connection-explicit}
\nabla^{\mathrm{sh}}_\cA
= d - \sum_{i < j} \omega_{ij}(\cA)\, d\log(z_i - z_j),
\end{equation}
where $\omega_{ij}(\cA) = \Res^{\mathrm{coll}}_{0,2}(\Theta_\cA)$
is the collision residue (the classical $r$-matrix).
\end{definition}

\begin{theorem}[Shadow/KZ duality on the affine comparison surface; \ClaimStatusProvedHere]
\label{thm:thqg-IV-kz-duality}
\index{KZ connection!duality}
The Verdier involution exchanges the shadow connections:
\begin{equation}\label{eq:thqg-IV-kz-duality}
\sigma(\nabla^{\mathrm{sh}}_\cA) = \nabla^{\mathrm{sh}}_{\cA^!}.
\end{equation}
For affine $\widehat{\fg}_k$, on the affine Kac--Moody comparison
surface this corresponds to the Feigin--Frenkel involution
$k \mapsto -k - 2h^\vee$ on the KZ connection
\[
\nabla_{\mathrm{KZ}}^{(k)}
= d - \frac{1}{k + h^\vee}
\sum_{i < j} \Omega_{ij}\, d\log(z_i - z_j),
\]
where $\Omega_{ij} = \sum_a t^a_i \otimes t^a_j$ is the Casimir.
\end{theorem}

\begin{proof}
The shadow connection is a genus-$0$ projection of $\Theta_\cA$.
By Theorem~\ref{thm:thqg-IV-theta-duality}, $\sigma(\Theta_\cA) = \Theta_{\cA^!}$.
Projecting to genus~$0$, arity~$n$:
$\sigma(\Sh_{0,n}(\Theta_\cA)) = \Sh_{0,n}(\Theta_{\cA^!})$.

For affine $\widehat{\fg}_k$, on the affine Kac--Moody comparison
surface the shadow/KZ comparison identifies the genus-$0$ shadow
connection with the KZ connection with coupling $1/(k + h^\vee)$.  Under
$k \mapsto k' = -k - 2h^\vee$:
$1/(k' + h^\vee) = 1/(-k - h^\vee) = -1/(k + h^\vee)$.
Therefore
$\nabla_{\mathrm{KZ}}^{(k')} = d + \frac{1}{k+h^\vee}
\sum_{i<j}\Omega_{ij}\,d\log(z_i - z_j)$,
which is the affine comparison-surface KZ connection with negated
coupling, i.e.\ the Feigin--Frenkel dual level.
\end{proof}

\subsubsection{BPZ duality}

\begin{proposition}[BPZ duality from $S$-duality; \ClaimStatusProvedHere]
\label{prop:thqg-IV-bpz-duality}
\index{BPZ duality!from $S$-duality}
For the Virasoro algebra, the duality $c \mapsto 26 - c$ sends
the $n$-point BPZ differential equation to its dual:
\begin{equation}\label{eq:thqg-IV-bpz-exchange}
\nabla^{\mathrm{BPZ}}_c \;\longmapsto\; \nabla^{\mathrm{BPZ}}_{26-c}.
\end{equation}
The BPZ parameter $b^2 = 6/c$ transforms to
$(b')^2 = 6/(26-c)$, and the Liouville parametrization
$c = 1 + 6(b + b^{-1})^2$ transforms by $b \mapsto b'$.
\end{proposition}

\begin{proof}
The Virasoro shadow connection at $n$ points acts on the conformal
block space $\cV_{c, \{h_i\}}(z_1, \ldots, z_n)$.
For a degenerate insertion at $h = h_{2,1} = (2b^2 - 1)/2$,
the null vector gives the BPZ equation:
\[
\left[
  \frac{b^2}{2}\,\partial_{z_0}^2
  + \sum_{i=1}^n \left(
    \frac{h_i}{(z_0 - z_i)^2}
    + \frac{\partial_{z_i}}{z_0 - z_i}
  \right)
\right]
\mathcal{F} = 0.
\]
Under $c \mapsto 26 - c$: the conformal weights $h_i$ transform
according to the Koszul dual module structure, and $b^2 \mapsto 6/(26-c)$.
The structural form of the equation is preserved because it is
determined by the shadow connection, which transforms via
$\sigma$ by Theorem~\ref{thm:thqg-IV-kz-duality}.
\end{proof}

\begin{remark}[BPZ parameter as Koszul coordinate]%
\label{rem:bpz-koszul-coordinate}%
\index{BPZ duality!Koszul coordinate}%
The BPZ parameter $b^2 = 6/c$ is a coordinate on the moduli of
Virasoro algebras.  Under Koszul duality $c \mapsto 26 - c$,
it transforms to $(b')^2 = 6/(26-c)$.  The product
\[
b^2 \cdot (b')^2 \;=\; \frac{36}{c(26-c)}
\]
is symmetric under $c \mapsto 26 - c$.  The Liouville
parametrization $c = 1 + 6(b + b^{-1})^2$ is the coordinate on
the self-dual family: the self-dual point $c = 13$ (where
$\mathrm{Vir}_c \cong \mathrm{Vir}_{26-c}$) corresponds to
$b^2 = 6/13$.
\end{remark}

\begin{proposition}[Shadow connection $S$-duality;
\ClaimStatusProvedHere]%
\label{prop:shadow-connection-s-duality}%
\index{shadow connection!$S$-duality}%
The shadow connection
$\nabla^{\mathrm{sh}} = d - \frac{Q'}{2Q}\,dt$
\textup{(}Theorem~\textup{\ref{thm:shadow-connection}}\textup{)}
transforms under $c \mapsto 26 - c$ by
\begin{equation}\label{eq:thqg-IV-shadow-connection-dual}
\nabla^{\mathrm{sh}}_{26-c}
\;=\;
d - \frac{(Q_{26-c})'}{2\,Q_{26-c}}\,dt.
\end{equation}
The complementarity of discriminants
$\Delta(c) + \Delta(26-c) = \mathrm{const}$
\textup{(}Proposition~\textup{\ref{prop:thqg-IV-quartic-duality}}\textup{)}
implies that the monodromy representations of
$\nabla^{\mathrm{sh}}_c$ and $\nabla^{\mathrm{sh}}_{26-c}$ are
related by the Koszul involution on the local system.
At the self-dual point $c = 13$, the involution $c \mapsto 26-c$
is the identity, so $Q_{\mathrm{Vir}_{13}}$ is invariant and
$\nabla^{\mathrm{sh}}_{13}$ is self-dual.
\end{proposition}

\begin{proof}
The shadow metric transforms as
$Q_{\mathrm{Vir}_c}(t) \mapsto Q_{\mathrm{Vir}_{26-c}}(t)$
under $c \mapsto 26 - c$, since Koszul duality replaces
$\kappa(c)$ by $\kappa(26-c) = 13 - \kappa(c)$
(Proposition~\ref{prop:thqg-IV-kappa-duality}) and the
quartic data by its dual
(Proposition~\ref{prop:thqg-IV-quartic-duality}).
The shadow connection is determined by the shadow metric via
$\nabla^{\mathrm{sh}} = d - \frac{Q'}{2Q}\,dt$; applying
the substitution gives~\eqref{eq:thqg-IV-shadow-connection-dual}.
The monodromy of $\nabla^{\mathrm{sh}}_c$ at a zero of $Q_c$ has
residue $1/2$, hence monodromy $-1$.
The shadow metric $Q_c(t) = (2\kappa(c) + 3\alpha(c)t)^2 + 2\Delta(c)t^2$ and its dual $Q_{26-c}(t) = (2\kappa(26-c) + 3\alpha(26-c)t)^2 + 2\Delta(26-c)t^2$ have zeros exchanged by the involution: a simple zero $t_0$ of $Q_c$ with $\mathrm{Res}_{t_0}(Q'/2Q) = 1/2$ maps to a zero of $Q_{26-c}$ with the same residue.  Since the monodromy around a pole of residue $1/2$ is $e^{2\pi i/2} = -1$, the monodromy representations of $\nabla^{\mathrm{sh}}_c$ and $\nabla^{\mathrm{sh}}_{26-c}$ are intertwined by the Koszul involution $\sigma\colon \mathrm{Vir}_c \to \mathrm{Vir}_{26-c}$.
At $c = 13$, $\kappa(13) = 13/2$ and $26-c = c$, so $Q_{\mathrm{Vir}_{13}}$ is invariant under $\sigma$.
\end{proof}

\begin{remark}[Functional equation of the shadow generating function]%
\label{rem:functional-equation-shadow-gf}%
\index{shadow generating function!functional equation}%
The $S$-duality involution $c \mapsto 26 - c$ induces a functional
equation on the shadow generating function.
At the arity-$2$ level,
$\kappa(c) + \kappa(26-c) = 13$ (constant)
(Proposition~\ref{prop:thqg-IV-kappa-duality}).
At arity~$4$, the discriminant complementarity
\[
\Delta(c) + \Delta(26-c)
\;=\;
\frac{6960}{(5c+22)(152-5c)}
\]
(Proposition~\ref{prop:thqg-IV-quartic-duality})
provides the quartic functional equation.
This is the shadow-level manifestation of string-theoretic
modular invariance: the Koszul involution on the Virasoro family
acts as a functional equation on the shadow invariants at each
arity, with the constant $13$ at arity~$2$ reflecting the critical
dimension $d = 26$ of the bosonic string.
\end{remark}

\begin{remark}[Koszul duality as functional equation]
\label{rem:koszul-duality-functional-equation}
\index{Koszul duality!functional equation interpretation}
The Virasoro shadow generating function $H(t,c)$ at central charge~$c$ and its Koszul dual $H(t, 26-c)$ at central charge~$26-c$ are related by the involution $c \mapsto 26-c$.  At arity~$2$: $\kappa(c) + \kappa(26-c) = 13$ (constant).  At arity~$4$: $\Delta(c) + \Delta(26-c) = 6960/[(5c+22)(152-5c)]$ (Proposition~\ref{prop:thqg-IV-quartic-duality}).  At all arities: the sigma-invariant $\Delta^{(r)} := S_r(c) + S_r(26-c)$ is level-independent (i.e., independent of $c$) if and only if $r \leq 3$.  For $r \geq 4$, the sigma-invariant is a nontrivial rational function of~$c$ whose denominator contains $c^{r-4}(5c+22)^{\lfloor(r-4)/2\rfloor}$: the dynamical content begins at arity~$4$.
\end{remark}

\subsubsection{Derived equivalence of local systems}

\begin{theorem}[Local system duality; \ClaimStatusProvedHere]
\label{thm:thqg-IV-local-system-duality}
\index{local systems!Koszul duality}
The Verdier involution induces an equivalence of derived categories
of local systems:
\begin{equation}\label{eq:thqg-IV-loc-equivalence}
\sigma^*\colon
D^b(\Loc_{\cA}(\Conf_n(X)))
\;\xrightarrow{\;\sim\;}
D^b(\Loc_{\cA^!}(\Conf_n(X))),
\end{equation}
where $\Loc_{\cA}$ denotes the category of local systems with
fiber~$V$ and monodromy determined by the shadow connection
$\nabla^{\mathrm{sh}}_\cA$.
\end{theorem}

\begin{proof}
The shadow connection $\nabla^{\mathrm{sh}}_\cA$ is a flat connection on
$\Conf_n(X)$ (flatness $(\nabla^{\mathrm{sh}})^2 = 0$ is
the genus-$0$ MC equation for $\Theta_\cA$, which holds by
Theorem~\ref{thm:mc2-bar-intrinsic}).  Its monodromy defines a
representation $\rho_\cA\colon \pi_1(\Conf_n(X)) \to \mathrm{GL}(V)$.

The functor $\sigma^*$ sends a local system $(L, \nabla)$ with
connection form $A \in \Omega^1(\Conf_n) \otimes \End(V)$ to the
local system $(L', \nabla')$ with $A' = \sigma(A)$.  Since $\sigma$
is a dg~Lie isomorphism preserving flatness
(Proposition~\ref{prop:thqg-IV-verdier-dg-lie}), $\sigma^*$ preserves
flatness and defines an exact functor on derived categories.
The inverse $(\sigma^{-1})^*$ gives the quasi-inverse.
\end{proof}

\subsubsection{Monodromy duality}

\begin{proposition}[Monodromy eigenvalue reciprocation; \ClaimStatusProvedHere]
\label{prop:thqg-IV-monodromy-duality}
\index{monodromy!duality}
For the KZ connection on $\widehat{\fg}_k$-modules, the braiding
eigenvalues $\lambda_i$ of the monodromy representation at level~$k$
and $\lambda'_i$ at level $k' = -k - 2h^\vee$ satisfy
\begin{equation}\label{eq:thqg-IV-monodromy-reciprocation}
\lambda'_i = \lambda_i^{-1}.
\end{equation}
\end{proposition}

\begin{proof}
The braiding eigenvalue for the exchange of two highest-weight modules
$V_\mu, V_\nu$ in the fusion channel $V_\rho$ is
\[
\lambda_{\mu,\nu}^\rho
= e^{2\pi i (h_\rho - h_\mu - h_\nu)/(k + h^\vee)},
\]
where $h_\mu = (\mu, \mu + 2\rho_{\mathrm{Weyl}})/(2(k + h^\vee))$ is
the conformal weight.  Under $k \mapsto k' = -k - 2h^\vee$:
$k' + h^\vee = -(k + h^\vee)$, so the conformal weights are negated
and $\lambda' = e^{-2\pi i(\cdots)} = \lambda^{-1}$.
\end{proof}

\begin{remark}[Virasoro braiding under duality]
\label{rem:thqg-IV-virasoro-braiding}
\index{braiding!Virasoro duality}
For the Virasoro algebra, the braiding eigenvalues of the degenerate
module $V_{h_{2,1}}$ are $\sigma_\pm = e^{\pm i\pi b^{-1}\alpha_2}$
(equation~\eqref{eq:degenerate-eigenvalues}).  Under $c \mapsto 26-c$:
$b^2 = 6/c \mapsto 6/(26-c)$, and the eigenvalues transform by a
nontrivial algebraic map.  At the self-dual point $c = 13$:
$b^2 = 6/13$ and the braiding is self-conjugate.
\end{remark}

% ======================================================================
%
%  IV. THE COMPLEMENTARITY CONSTANT AND THE EXPONENT SUM FORMULA
%
% ======================================================================

\subsection{Complementarity constants for the standard landscape}
\label{subsec:thqg-IV-complementarity-constant}
\index{complementarity!constant $K(\cA)$|textbf}

\begin{definition}[Complementarity constant]
\label{def:thqg-IV-complementarity-constant}
\index{complementarity!constant|textbf}
For a chiral Koszul pair $(\cA, \cA^!)$, the
\emph{complementarity constant} is
\begin{equation}\label{eq:thqg-IV-K-def}
K(\cA)
\;:=\;
\kappa(\cA) + \kappa(\cA^!).
\end{equation}
By Theorem~\ref{thm:modular-characteristic}\textup{(}iii\textup{)},
$K(\cA)$ is the trace of
$\Theta_\cA + \sigma(\Theta_\cA)$ at arity~$2$:
\[
K(\cA)
= \Tr\bigl(\Theta_\cA^{(1)} + \Theta_{\cA^!}^{(1)}\bigr)
= \Tr\bigl(\Theta_\cA^{(1)} + \sigma(\Theta_\cA^{(1)})\bigr),
\]
where $\Theta^{(1)}$ denotes the genus-$1$ component.
\end{definition}

\begin{remark}[Physical interpretation]
\label{rem:thqg-IV-K-physical}
The complementarity constant $K(\cA)$ is the total genus-$1$
curvature of the Koszul pair: it measures the combined failure of the
bar differentials of~$\cA$ and~$\cA^!$ to square to zero at
genus~$1$.  When $K(\cA) = 0$, the curvatures cancel exactly and the
pair is \emph{genus-$1$-flat}.  When $K(\cA) \neq 0$, there is a
residual curvature shared by the pair, originating in the nonlinearity
of the Drinfeld--Sokolov parametrization.
\end{remark}

\subsubsection{Heisenberg algebra}
\index{Heisenberg algebra!complementarity constant}

\begin{proposition}[Heisenberg complementarity; \ClaimStatusProvedHere]
\label{prop:thqg-IV-heisenberg-K}
For the Heisenberg algebra $\cH_\kappa$ at level $\kappa \neq 0$:
\begin{equation}\label{eq:thqg-IV-heisenberg-K}
\kappa(\cH_\kappa) = \kappa,
\qquad
\kappa(\cH_\kappa^!) = -\kappa,
\qquad
K(\cH_\kappa) = 0.
\end{equation}
\end{proposition}

\begin{proof}
The Heisenberg OPE is $J(z) J(w) \sim \kappa/(z-w)^2$.  The bar
differential is determined by the double-pole coefficient alone,
giving $\kappa(\cH_\kappa) = \kappa$
(Theorem~\ref{thm:heisenberg-higher-genus}).
The Koszul dual is $\cH_\kappa^! = \operatorname{Sym}^{\mathrm{ch}}(V^*)$,
with $\kappa(\cH_\kappa^!) = -\kappa$: Verdier duality
on $\overline{C}_2(X)$ negates the propagator $\eta_{12}$, hence
the curvature sign.  The genus-$1$ curvature is $d_{\mathrm{fib}}^2 = \kappa \omega_1$;
under duality, $\kappa \mapsto -\kappa$ and $\omega_1$ is unchanged
(it is a topological datum).  Therefore $K = \kappa + (-\kappa) = 0$.
\end{proof}

\subsubsection{Affine Kac--Moody algebras}
\index{Kac--Moody algebra!complementarity constant}

\begin{theorem}[Affine Kac--Moody complementarity; \ClaimStatusProvedHere]
\label{thm:thqg-IV-km-K}
For any simple Lie algebra $\fg$ with dual Coxeter number $h^\vee$
and dimension $d = \dim\fg$, the affine Kac--Moody algebra
$\widehat{\fg}_k$ at level $k \neq -h^\vee$ satisfies:
\begin{equation}\label{eq:thqg-IV-km-kappa}
\kappa(\widehat{\fg}_k)
\;=\;
\frac{(k + h^\vee)\, d}{2h^\vee},
\end{equation}
the Koszul dual is
$\widehat{\fg}_k^! = \mathrm{CE}^{\mathrm{ch}}(\widehat{\fg}_{k'})$
with $k' = -k - 2h^\vee$
(same~$\kappa$ as $\widehat{\fg}_{k'}$, categorically distinct), and
\begin{equation}\label{eq:thqg-IV-km-K}
K(\widehat{\fg}_k) = 0.
\end{equation}
\end{theorem}

\begin{proof}
The genus-$1$ curvature receives two contributions from the OPE
$J^a(z)J^b(w) \sim k\kappa^{ab}/(z-w)^2 + f^{abc}J_c(w)/(z-w)$.

\emph{Double-pole channel.}
The double pole $k\kappa^{ab}/(z-w)^2$ contributes through the
self-loop graph: one vertex of genus~$0$ with two legs and one
self-edge.  The contracted amplitude is
$\sum_{a,b} k\kappa^{ab} \cdot \kappa_{ab} = k \cdot \dim\fg$.
The Hodge normalization on $\overline{\cM}_{1,1}$ contributes a factor
of $1/(2h^\vee)$.  Net contribution: $kd/(2h^\vee)$.

\emph{Simple-pole channel.}
The simple pole $f^{abc}J_c(w)/(z-w)$ contributes through the tree
graph with two vertices of genus~$0$ connected by one edge.  The
contracted amplitude involves
$\sum_{a,b,c} f^{abc}f_{abc} = 2h^\vee \cdot \dim\fg$.
With the combinatorial factor and Hodge normalization: $d/2$.

Combining: $\kappa(\widehat{\fg}_k)
= kd/(2h^\vee) + d/2 = (k + h^\vee)d/(2h^\vee)$.

Under $k \mapsto k' = -k - 2h^\vee$:
$k' + h^\vee = -k - h^\vee$, so
$\kappa(\widehat{\fg}_{k'})
= (-k - h^\vee)d/(2h^\vee) = -(k + h^\vee)d/(2h^\vee) = -\kappa(\widehat{\fg}_k)$.
Therefore $K = 0$.
\end{proof}

\begin{computation}[Affine KM complementarity for classical types]
\label{comp:thqg-IV-km-classical}
\index{complementarity!classical Lie algebras}
\renewcommand{\arraystretch}{1.4}
\begin{center}
\begin{tabular}{lcccccc}
\toprule
$\fg$ & $d$ & $h^\vee$ & $\kappa(\widehat{\fg}_k)$ & $k'$ &
$\kappa(\widehat{\fg}_{k'})$ & $K$ \\
\midrule
$\fsl_2$ & $3$ & $2$ & $\dfrac{3(k+2)}{4}$ & $-k-4$ &
$\dfrac{-3(k+2)}{4}$ & $0$ \\[6pt]
$\fsl_3$ & $8$ & $3$ & $\dfrac{4(k+3)}{3}$ & $-k-6$ &
$\dfrac{-4(k+3)}{3}$ & $0$ \\[6pt]
$\fsl_N$ & $N^2{-}1$ & $N$ & $\dfrac{(k{+}N)(N^2{-}1)}{2N}$ & $-k-2N$ &
$\dfrac{-(k{+}N)(N^2{-}1)}{2N}$ & $0$ \\[6pt]
$\fso_{2N}$ & $N(2N{-}1)$ & $2N{-}2$ & $\dfrac{N(2N{-}1)(k{+}2N{-}2)}{2(2N{-}2)}$ & $-k-4N+4$ &
$-\kappa$ & $0$ \\[6pt]
$E_8$ & $248$ & $30$ & $\dfrac{124(k+30)}{30}$ & $-k-60$ &
$\dfrac{-124(k+30)}{30}$ & $0$ \\
\bottomrule
\end{tabular}
\end{center}
The exact cancellation $K = 0$ is independent of the Lie type: it depends
only on the linearity of $\kappa$ in the shifted level $t = k + h^\vee$.
\end{computation}

\subsubsection{\texorpdfstring{$bc$--$\beta\gamma$}{bc-betagamma} systems}
\index{bc system@$bc$ system!complementarity constant}

\begin{proposition}[\texorpdfstring{$bc$--$\beta\gamma$}{bc-betagamma}
complementarity; \ClaimStatusProvedHere]
\label{prop:thqg-IV-bc-bg-K}
For the $bc$ ghost system of conformal weight $\lambda$:
\begin{equation}\label{eq:thqg-IV-bc-bg-K}
\kappa(bc_\lambda) = -(6\lambda^2 - 6\lambda + 1),
\qquad
\kappa(\beta\gamma_\lambda) = 6\lambda^2 - 6\lambda + 1,
\qquad
K = 0.
\end{equation}
\end{proposition}

\begin{proof}
The central charge of $bc_\lambda$ is $c = -2(6\lambda^2 - 6\lambda + 1)$.
The $bc$ system has a single generator pair $(b, c)$ of conformal weights
$(\lambda, 1 - \lambda)$; the double-pole coefficient is $1/(z-w)^1$
(first-order pole only for $\lambda \neq 1$), but the genus-$1$
curvature receives a contribution from the period matrix of the
fermion determinant, giving $\kappa = c/2 = -(6\lambda^2 - 6\lambda + 1)$.

The Koszul dual of $bc_\lambda$ is $\beta\gamma_\lambda$: the
$\mathrm{Ext}/\mathrm{Sym}$ duality (exterior $\leftrightarrow$
symmetric) negates the central charge, giving
$c(\beta\gamma_\lambda) = 2(6\lambda^2 - 6\lambda + 1)$ and
$\kappa(\beta\gamma_\lambda) = 6\lambda^2 - 6\lambda + 1$.
Therefore $K = 0$.

At the standard ghost values $\lambda = 2$: $\kappa(bc_2) = -13$,
$\kappa(\beta\gamma_2) = 13$.
At $\lambda = 1$: $\kappa(bc_1) = -1$, $\kappa(\beta\gamma_1) = 1$.
\end{proof}

\subsubsection{Lattice vertex algebras}

\begin{proposition}[Lattice complementarity; \ClaimStatusProvedHere]
\label{prop:thqg-IV-lattice-K}
For a lattice VOA $V_\Lambda$ of rank~$r$:
$\kappa(V_\Lambda) = r$, $\kappa(V_\Lambda^!) = -r$, $K = 0$.
\end{proposition}

\begin{proof}
By additivity of $\kappa$ (Corollary~\ref{cor:kappa-additivity}) and
independence of the lattice cocycle
(Theorem~\ref{thm:lattice:curvature-braiding-orthogonal}):
$\kappa(V_\Lambda) = r \cdot \kappa(\cH_1) = r$.
The Koszul dual negates each Heisenberg factor: $\kappa(V_\Lambda^!) = -r$.
\end{proof}

\subsubsection{Virasoro algebra}
\index{Virasoro algebra!complementarity constant}

\begin{theorem}[Virasoro complementarity; \ClaimStatusProvedHere]
\label{thm:thqg-IV-virasoro-K}
For $\mathrm{Vir}_c$:
\begin{equation}\label{eq:thqg-IV-virasoro-kappa}
\kappa(\mathrm{Vir}_c) = \frac{c}{2},
\qquad
\mathrm{Vir}_c^! = \mathrm{Vir}_{26-c},
\qquad
K(\mathrm{Vir}_c) = 13.
\end{equation}
\end{theorem}

\begin{proof}
The Virasoro OPE $T(z) T(w) \sim (c/2)/(z-w)^4 + 2T(w)/(z-w)^2
+ \partial T(w)/(z-w)$ gives a quartic pole with coefficient $c/2$.
The genus-$1$ curvature is $\kappa(\mathrm{Vir}_c) = c/2$
(Theorem~\ref{thm:vir-genus1-curvature}).

The Koszul dual central charge: the Virasoro algebra is obtained from
$\widehat{\fsl}_2$ by quantum Drinfeld--Sokolov reduction.
The DS central charge formula is
$c(k) = 1 - 6(k+2)^{-1}(k+2-1)^2/(k+2)
= 13 - 6(k+2) - 6/(k+2)$,
with $c + c' = 26$ under $k \mapsto k' = -k - 4$
(Theorem~\ref{thm:central-charge-complementarity}).
Therefore
$K = c/2 + (26-c)/2 = 13$.

The nonzero $K$ arises because the DS parametrization
$c(k) = 1 - 6(k+1)^2/(k+2)$ is quadratic in~$k$: the Feigin--Frenkel
involution sends $c \mapsto 26 - c$, not $c \mapsto -c$.  Hence
$\kappa' = (26-c)/2 \neq -\kappa$.
\end{proof}

\begin{remark}[The Virasoro anomaly]
\label{rem:thqg-IV-virasoro-anomaly}
\index{Virasoro!anomaly origin of $K \neq 0$}
The nonvanishing of $K(\mathrm{Vir}_c) = 13$ has a direct
representation-theoretic origin: the conformal weight of the stress tensor
$T$ is $h = 2$, and the ghost system has $c_{\mathrm{ghost}} = -26$.
The constant $13 = 26/2$ is the effective curvature of $26$ free bosons
after ghost subtraction.  Equivalently, $13$ is the value of $\kappa$
at the self-dual point $c = 13$, and by the functional equation
(Theorem~\ref{thm:thqg-IV-free-energy}) it is the level-independent
residue of the Koszul pair.
\end{remark}

\subsubsection{Principal \texorpdfstring{$\mathcal{W}$}{W}-algebras}
\index{W-algebra@$\mathcal{W}$-algebra!complementarity constant}

\begin{theorem}[Principal $\mathcal{W}$-algebra complementarity; \ClaimStatusProvedHere]
\label{thm:thqg-IV-w-K}
For $\mathcal{W}_N := \mathcal{W}^k(\fsl_N, f_{\mathrm{prin}})$:
\begin{equation}\label{eq:thqg-IV-w-kappa}
\kappa(\mathcal{W}_N^k) = c(k) \cdot (H_N - 1),
\end{equation}
where $H_N = \sum_{j=1}^N 1/j$ is the $N$-th harmonic number.  The complementarity constant is
\begin{equation}\label{eq:thqg-IV-w-K}
K(\mathcal{W}_N) = K_N \cdot (H_N - 1),
\qquad K_N := c + c' = 2(N{-}1)(2N^2{+}2N{+}1).
\end{equation}
\end{theorem}

\begin{proof}
The $\mathcal{W}_N$ algebra has generators $W_2 = T, W_3, \ldots, W_N$
of conformal weights $2, 3, \ldots, N$.  The exponents of $\fsl_N$ are
$m_i = i$ for $i = 1, \ldots, N-1$.  The exponent-sum invariant:
\[
\varrho(\fsl_N)
= \sum_{i=1}^{N-1} \frac{1}{m_i + 1}
= \sum_{i=1}^{N-1} \frac{1}{i+1}
= \sum_{s=2}^N \frac{1}{s}
= H_N - 1.
\]
By the anomaly ratio (Corollary~\ref{cor:anomaly-ratio}):
$\kappa(\mathcal{W}_N^k) = c(k) \cdot \varrho(\fsl_N) = c(k)(H_N - 1)$.

The DS central charge for $\fsl_N$ is
$c(k) = (N-1) - N(N^2-1)(k+N-1)^2/(k+N)$;
equivalently, writing $t = k+N$: $c = (N-1)(2N^2+2N+1) - N(N^2-1)(t + 1/t)$.
The Feigin--Frenkel involution $k \mapsto k' = -k - 2N$ sends $t \mapsto -t$,
so $c' = (N-1)(2N^2+2N+1) - N(N^2-1)(-t - 1/t)
= (N-1)(2N^2+2N+1) + N(N^2-1)(t + 1/t)$.  Therefore
$c + c' = 2(N-1)(2N^2+2N+1) = K_N$
(Theorem~\ref{thm:central-charge-complementarity}),
independent of~$k$.
Therefore $K = K_N \cdot (H_N - 1)$.
\end{proof}

\begin{computation}[Complementarity constants for small $N$]
\label{comp:thqg-IV-w-small-N}
\renewcommand{\arraystretch}{1.4}
\begin{center}
\begin{tabular}{cccccc}
\toprule
$N$ & $H_N - 1$ & $K_N = c + c'$ & $K(\mathcal{W}_N)$ &
Self-dual $c_*$ & $\kappa(c_*)$ \\
\midrule
$2$ & $1/2$ & $26$ & $13$ & $13$ & $13/2$ \\[4pt]
$3$ & $5/6$ & $100$ & $250/3$ & $50$ & $250/6$ \\[4pt]
$4$ & $13/12$ & $246$ & $1599/6$ & $123$ & $1599/12$ \\[4pt]
$5$ & $77/60$ & $488$ & $37576/60$ & $244$ & $37576/120$ \\[4pt]
$6$ & $29/20$ & $850$ & $24650/20$ & $425$ & $24650/40$ \\
\bottomrule
\end{tabular}
\end{center}
The self-dual central charge is $c_* = K_N/2 = (N-1)(2N^2+2N+1)$.
The complementarity constant $K(\mathcal{W}_N)$ grows as
$\sim 4N^3 \log N / 3$ for large $N$, reflecting the logarithmic
growth of harmonic numbers.
\end{computation}

\subsubsection{The Koszul conductor and its growth}

\begin{definition}[Koszul conductor]
\label{def:thqg-IV-koszul-conductor}
\index{Koszul conductor|textbf}
The \emph{Koszul conductor} is the level-independent central charge sum
\begin{equation}\label{eq:thqg-IV-koszul-conductor}
K_N := c(\mathcal{W}_N^k) + c(\mathcal{W}_N^{k'})
= 2(N-1)(2N^2 + 2N + 1).
\end{equation}
\end{definition}

\begin{proposition}[Koszul conductor asymptotics; \ClaimStatusProvedHere]
\label{prop:thqg-IV-conductor-growth}
\index{Koszul conductor!asymptotics}
For large~$N$:
\begin{equation}\label{eq:thqg-IV-conductor-growth}
K_N = 4N^3 + O(N^2),
\qquad
K(\mathcal{W}_N) = 4N^3 \log N + O(N^3).
\end{equation}
\end{proposition}

\begin{proof}
$K_N = 2(N-1)(2N^2+2N+1) = 4N^3 - 2N - 2$.
$H_N - 1 = \log N + \gamma - 1 + O(1/N)$, so
$K(\mathcal{W}_N) = K_N(H_N - 1) = 4N^3 \log N + O(N^3)$.
\end{proof}

\subsubsection{General simple Lie algebras: the exponent sum formula}

\begin{theorem}[Exponent sum formula; \ClaimStatusProvedHere]
\label{thm:thqg-IV-exponent-sum}
\index{exponent sum formula|textbf}
For any simple $\fg$ of rank~$r$ with exponents $m_1, \ldots, m_r$:
\begin{equation}\label{eq:thqg-IV-exponent-sum}
K\bigl(\mathcal{W}(\fg)\bigr)
= (c_{\fg} + c_{\fg}') \cdot \varrho(\fg),
\qquad
\varrho(\fg) := \sum_{i=1}^r \frac{1}{m_i + 1}.
\end{equation}
For affine KM: $K(\widehat{\fg}_k) = 0$.
\end{theorem}

\begin{proof}
By the anomaly ratio $\kappa = c \cdot \varrho(\fg)$
(Corollary~\ref{cor:anomaly-ratio}): $\kappa' = c' \cdot \varrho(\fg)$.
The central charge sum $c + c'$ is level-independent
(Theorem~\ref{thm:central-charge-complementarity}).
Therefore $K = \kappa + \kappa' = (c + c') \cdot \varrho = K_{\fg} \cdot \varrho$.

For affine KM, $\kappa = (k+h^\vee)d/(2h^\vee)$ is linear in
$k + h^\vee$; the involution $k \mapsto -k - 2h^\vee$ sends
$k + h^\vee \mapsto -(k + h^\vee)$, giving $K = 0$.
\end{proof}

\begin{computation}[Exponent sums for all simple types]
\label{comp:thqg-IV-exponent-sums}
\index{complementarity!all simple types}
\renewcommand{\arraystretch}{1.4}
\begin{center}
\begin{tabular}{lccccc}
\toprule
$\fg$ & Exponents & $\varrho(\fg)$ & $c_{\mathrm{KM}}+c_{\mathrm{KM}}'$ &
$c_{\mathcal{W}}+c_{\mathcal{W}}'$ & $K(\mathcal{W}(\fg))$ \\
\midrule
$A_1$ & $1$ & $1/2$ & $6$ & $26$ & $13$ \\
$A_2$ & $1, 2$ & $5/6$ & $16$ & $100$ & $250/3$ \\
$A_3$ & $1, 2, 3$ & $13/12$ & $30$ & $246$ & $1599/6$ \\
$A_4$ & $1, 2, 3, 4$ & $77/60$ & $48$ & $488$ & $37576/60$ \\
$B_2$ & $1, 3$ & $3/4$ & $20$ & $66$ & $99/2$ \\
$B_3$ & $1, 3, 5$ & $23/24$ & $42$ & $190$ & $4370/24$ \\
$D_4$ & $1, 3, 3, 5$ & $4/3$ & $56$ & $304$ & $1216/3$ \\
$G_2$ & $1, 5$ & $2/3$ & $28$ & $74$ & $148/3$ \\
$F_4$ & $1, 5, 7, 11$ & $7/8$ & $104$ & $640$ & $560$ \\
$E_6$ & $1, 4, 5, 7, 8, 11$ & $427/360$ & $156$ & $1248$ & $532896/360$ \\
$E_7$ & $1, 5, 7, 9, 11, 13, 17$ & $2777/2520$ & $266$ & $2674$ & $7424498/2520$ \\
$E_8$ & $1, 7, 11, 13, 17, 19, 23, 29$ & $121/126$ & $496$ & $5768$ & $697928/126$ \\
\bottomrule
\end{tabular}
\end{center}
Two features: $\varrho(\fg)$ is not monotone in rank
($\varrho(E_8) < \varrho(E_6)$, reflecting large exponents);
$K$ grows roughly as $d \cdot \log r$.
\end{computation}

\begin{remark}[Structural content of $K = 0$ vs.\ $K \neq 0$]
\label{rem:thqg-IV-K-zero-vs-nonzero}
\index{complementarity!structural interpretation}
The vanishing of $K$ for affine KM and the nonvanishing for
$\mathcal{W}$-algebras has a single root cause:
\emph{the Drinfeld--Sokolov reduction is nonlinear}.
For affine KM, $\kappa$ is linear in $k + h^\vee$; the
Feigin--Frenkel involution negates it.
For $\mathcal{W}$-algebras, $\kappa = c(k) \cdot \varrho$ where
$c(k)$ is rational in~$k$; the involution does not negate~$\kappa$.
The nonlinearity of DS reduction is therefore detected at the level
of complementarity constants, and $K \neq 0$ is the hallmark of a
genuinely non-linear chiral algebra.
\end{remark}

\begin{computation}['t~Hooft coupling under complementarity]
\label{comp:thqg-IV-thooft-complementarity}
\index{'t Hooft coupling!complementarity constant}
Write $\lambda = N/(k+N)$.
The Feigin--Frenkel involution sends
$k + N \mapsto -(k+N)$, hence $\lambda \mapsto -\lambda$.
For affine $\widehat{\fsl}_{N,k}$:
\[
c + c'
= \frac{kd}{k+N} + \frac{(-k-2N)d}{-k-N}
= \frac{kd + (k+2N)d}{k+N}
= 2d,
\qquad
K = 0.
\]
For $\mathcal{W}_N^k$, the DS formula
$c(k) = (N{-}1) - N(N^2{-}1)(k{+}N{-}1)^2/(k{+}N)$
is quadratic in~$k$; the central charge sum retains the residue
$K_N = 2(N{-}1)(2N^2{+}2N{+}1)$, giving $K \neq 0$.
The Gaberdiel--Gopakumar involution
$\lambda \mapsto 1 - \lambda$ is the level-rank involution
$(N,k) \leftrightarrow (k,N)$, distinct from the
Feigin--Frenkel involution $\lambda \mapsto -\lambda$.
\end{computation}

% ======================================================================
%
%  IV.5. THE DS COMPLEMENTARITY DEFECT
%
% ======================================================================

\subsubsection{The DS complementarity defect}
\label{subsec:thqg-IV-ds-defect}
\index{DS complementarity defect|textbf}
\index{Drinfeld--Sokolov reduction!complementarity defect}

\begin{definition}[DS complementarity defect]
\label{def:ds-complementarity-defect}
For a principal DS reduction
$\phi\colon \widehat{\fg}_k \to \mathcal{W}^k(\fg, f_{\mathrm{prin}})$,
the \emph{DS complementarity defect} is
\begin{equation}\label{eq:ds-defect}
\delta_K(\phi)
\;:=\;
K\bigl(\mathcal{W}^k(\fg)\bigr) - K(\widehat{\fg}_k).
\end{equation}
Since $K(\widehat{\fg}_k) = 0$
(Theorem~\ref{thm:thqg-IV-km-K}):
$\delta_K(\phi) = K(\mathcal{W}^k(\fg))
= K_\fg \cdot \varrho(\fg)$
where $K_\fg = 2r + 4dh^\vee$ is the Koszul conductor
and $\varrho(\fg) = \sum 1/(m_i + 1)$ is the exponent-sum
invariant.
\end{definition}

\begin{theorem}[DS complementarity tower; \ClaimStatusProvedHere]
\label{thm:ds-complementarity-tower}
\index{DS complementarity defect!tower}
Let $\fg$ be a simple Lie algebra with rank~$r$, dimension~$d$,
dual Coxeter number~$h^\vee$, and exponents $m_1, \ldots, m_r$.
Let $\phi\colon \widehat{\fg}_k \to \mathcal{W}(\fg)$ be
the principal DS reduction.
Define the \emph{DS complementarity tower} entry at arity~$s$
and genus~$g$:
\begin{equation}\label{eq:ds-tower-def}
\Delta_\phi^{(s,g)}
\;:=\;
\Sh_{s,g}\bigl(\mathcal{W}(\fg)\bigr)
\;+\;
\sigma\!\bigl(\Sh_{s,g}(\mathcal{W}(\fg))\bigr).
\end{equation}
\begin{enumerate}[label=\textup{(\roman*)}]
\item \emph{Affine vanishing.}
  $\Delta_{\mathrm{aff}}^{(s,g)} = 0$ at all $(s,g)$.

\item \emph{Arity $2$, all genera.}
  \begin{equation}\label{eq:ds-tower-arity2}
  \Delta_\phi^{(2,g)}
  \;=\;
  K_\fg \cdot \varrho(\fg) \cdot \lambda_g^{\mathrm{FP}},
  \qquad
  K_\fg = 2r + 4dh^\vee,
  \quad
  \varrho(\fg) = \textstyle\sum_{i=1}^r \frac{1}{m_i + 1},
  \end{equation}
  independent of the level~$k$.

\item \emph{Arity $3$, genus $0$.}
  The gravitational cubic has the same conformal weights for
  $\mathcal{W}(\fg)$ and~$\mathcal{W}(\fg)^!$:
  \begin{equation}\label{eq:ds-tower-cubic}
  \Delta_\phi^{(3,0)}
  \;=\;
  -K_\fg \cdot \varrho(\fg),
  \end{equation}
  independent of the level.
  For $\fg = \fsl_2$: $-K \cdot \varrho = -26 \cdot \tfrac{1}{2} = -13$;
  the full cubic conductor
  $\mathfrak{C}_c + \mathfrak{C}_{26-c} = -26$
  includes the Hessian normalization.

\item \emph{Arity $\geq 4$: level dependence.}
  For $\fg = \fsl_2$:
  \[
  \Delta_\phi^{(4,0)}
  \;=\;
  \frac{10\bigl(10(c-13)^2 + 2262\bigr)}
  {c(5c+22)(26-c)(152-5c)},
  \]
  which depends on~$c$.

\item \emph{Asymptotics \textup{(}type~$A$\textup{)}.}
  $\delta_K(\phi) \sim 4N^3 \log N$ for $N \to \infty$.
\end{enumerate}
\end{theorem}

\begin{proof}
\emph{(i)}
At arity~$2$: $\kappa = (k{+}h^\vee)d/(2h^\vee)$ is linear in the
shifted level $t = k + h^\vee$; the Feigin--Frenkel involution
$t \mapsto -t$ negates~$\kappa$, giving $\Delta^{(2,g)} = 0$
at every genus.
At arity~$3$: the affine OPE is associative
($m_3 = 0$; Computation~\ref{comp:thqg-IV-cubic-virasoro}), so
$\mathfrak{C}(\widehat{\fg}_k) = 0$ and
$\Delta^{(3,g)} = 0$.
At arity~$\geq 4$: the affine shadow tower terminates
($o^{(4)} = 0$ by the Jacobi identity), so $\Sh_s = 0$ for $s \geq 4$
and $\Delta^{(s,g)} = 0$.

\emph{(ii)}
$K(\mathcal{W}(\fg)) = K_\fg \cdot \varrho(\fg)$
(Theorem~\ref{thm:thqg-IV-w-K}), with
$K_\fg = c + c'$ level-independent.
The genus-$g$ entry is
$\Delta^{(2,g)} = K \cdot \lambda_g^{\mathrm{FP}}$
by scalar universality.

\emph{(iii)}
The cubic shadow on the $T$-line is
$\mathfrak{C} = -c \cdot \varrho$
(Sugawara normal-ordering in the spin-$2$ channel).
Under $\sigma$: $c \mapsto c' = K_\fg - c$.
Sum: $-(c + c') \cdot \varrho = -K_\fg \cdot \varrho$.

\emph{(iv)}
$Q^{\mathrm{contact}}_c = 10/(c(5c+22))$
and $Q^{\mathrm{contact}}_{K_2 - c} = 10/((26-c)(152-5c))$.
Numerator of the sum: $10c^2 - 260c + 3952 = 10(c-13)^2 + 2262$,
a quadratic in~$c$.

\emph{(v)}
$K_\fg = 4N^3 + O(N^2)$ for $\fg = \fsl_N$,
$\varrho = \log N + O(1)$.
\end{proof}

% ======================================================================
%
%  IV.6. THE SIGMA-INVARIANT SHADOW RING
%
% ======================================================================

\subsubsection{The $\sigma$-invariant shadow ring}
\label{subsec:thqg-IV-sigma-ring}
\index{sigma-invariant shadow ring@$\sigma$-invariant shadow ring|textbf}

The DS complementarity tower
$\{\Delta_\phi^{(s,g)}\}_{s,g}$ of
Theorem~\ref{thm:ds-complementarity-tower} is the bigraded piece
of a ring.

\begin{definition}[$\sigma$-invariant shadow ring]
\label{def:sigma-invariant-shadow-ring}
For a chiral Koszul pair $(\cA, \cA^!)$ with Verdier involution
$\sigma$, the \emph{$\sigma$-invariant shadow ring} is the
$\sigma$-fixed locus
\begin{equation}\label{eq:sigma-ring-def}
(\cA^{\mathrm{sh}})^\sigma
\;:=\;
\bigl\{x \in \cA^{\mathrm{sh}} : \sigma(x) = x\bigr\},
\end{equation}
a bigraded commutative subring of~$\cA^{\mathrm{sh}}$.
Its $(s,g)$-component is $\Delta^{(s,g)}
= \Sh_{s,g}(\cA) + \sigma(\Sh_{s,g}(\cA))$.
\end{definition}

\begin{corollary}[Level-independence filtration; \ClaimStatusProvedHere]
\label{cor:level-independence-filtration}
\index{level-independence filtration}
Theorem~\textup{\ref{thm:ds-complementarity-tower}} gives
$(\cA^{\mathrm{sh}})^\sigma$ a filtration by level independence:
\begin{equation}\label{eq:filtration}
\underbrace{
\Delta^{(2,\bullet)},\;\Delta^{(3,0)}
}_{\text{level-independent}}
\;\subset\;
(\cA^{\mathrm{sh}})^\sigma
\;\supset\;
\underbrace{
\Delta^{(4,\bullet)},\;\Delta^{(5,\bullet)},\;\ldots
}_{\text{level-dependent}}.
\end{equation}
The level-independent part $R^{\mathrm{univ}}$
is generated by $K_\fg \cdot \varrho(\fg)$
at arity~$2$ and $-K_\fg \cdot \varrho(\fg)$ at arity~$3$.
The level-dependent part $I^{\mathrm{dyn}}$
has first nonzero entry at arity~$4$.
The threshold is arity~$3$ for all simple~$\fg$.
\end{corollary}

\begin{proof}
Level independence at arities $2$--$3$: parts~(ii)--(iii) of
Theorem~\ref{thm:ds-complementarity-tower}; the argument uses only
that the affine OPE is linear in $k + h^\vee$, which holds for
every simple~$\fg$.
Level dependence at arity~$4$: part~(iv); the DS central charge
$c(k) = r - dh^\vee(t{-}1)^2/t$ is quadratic in~$t$, so
$Q(c) + Q(K_\fg - c)$ is not constant.
\end{proof}

\begin{computation}[Koszul conductor for all simple types]
\label{comp:koszul-conductor-all-types}
\index{Koszul conductor!all simple types}
\renewcommand{\arraystretch}{1.3}
\begin{center}
\small
\begin{tabular}{lcccccc}
\toprule
$\fg$ & $d$ & $h^\vee$ & $r$ & Exponents & $K_\fg$ & $\varrho(\fg)$ \\
\midrule
$A_1$ & $3$ & $2$ & $1$ & $(1)$ & $26$ & $1/2$ \\
$A_2$ & $8$ & $3$ & $2$ & $(1,2)$ & $100$ & $5/6$ \\
$A_3$ & $15$ & $4$ & $3$ & $(1,2,3)$ & $246$ & $13/12$ \\
$B_2$ & $10$ & $3$ & $2$ & $(1,3)$ & $124$ & $3/4$ \\
$G_2$ & $14$ & $4$ & $2$ & $(1,5)$ & $228$ & $2/3$ \\
$D_4$ & $28$ & $6$ & $4$ & $(1,3,3,5)$ & $680$ & $7/6$ \\
$F_4$ & $52$ & $9$ & $4$ & $(1,5,7,11)$ & $1880$ & $7/8$ \\
$E_6$ & $78$ & $12$ & $6$ & $(1,4,5,7,8,11)$ & $3756$ & $427/360$ \\
$E_8$ & $248$ & $30$ & $8$ & $(1,7,11,13,17,19,23,29)$ & $29776$ & $121/126$ \\
\bottomrule
\end{tabular}
\end{center}
The arity-$2$ universal entry at genus~$1$ is
$\Delta^{(2,1)} = K_\fg \cdot \varrho / 24$; for $E_8$:
$\Delta^{(2,1)} = 29776 \cdot \varrho(E_8) / 24$.
\end{computation}

\begin{theorem}[Recursive determination of $(\cA^{\mathrm{sh}})^\sigma$;
\ClaimStatusProvedHere]
\label{thm:sigma-ring-finite-generation}
\index{sigma-invariant shadow ring@$\sigma$-invariant shadow ring!finite generation}
For $\cA = \mathrm{Vir}_c$, the $\sigma$-invariant shadow sequence on the
primary line is recursively determined by three initial data:
\begin{equation}\label{eq:sigma-ring-generators}
\Delta^{(2)} = 13,
\qquad
\Delta^{(3)} = 4,
\qquad
\Delta^{(4)} = \frac{20(5c^2 - 130c + 1976)}
{c(c-26)(5c-152)(5c+22)}.
\end{equation}
The first two are level-independent
\textup{(}the perturbative generators\textup{)}.
The third depends on~$c$ \textup{(}the dynamical generator\textup{)}.
All higher entries are determined by the $\sigma$-invariant
master equation recursion:
\begin{align}
\Delta^{(5)} &= -\tfrac{1}{10}\,
  \bigl\{\Delta^{(3)},\, \Delta^{(4)}\bigr\}_H^\sigma,
\label{eq:sigma-ring-arity5} \\
\Delta^{(6)} &= -\tfrac{1}{12}\,
  \Bigl[\bigl\{\Delta^{(3)},\, \Delta^{(5)}\bigr\}_H^\sigma
  + \tfrac{1}{2}\bigl\{\Delta^{(4)},\, \Delta^{(4)}\bigr\}_H^\sigma
  \Bigr],
\label{eq:sigma-ring-arity6}
\end{align}
where $\{f, g\}_H^\sigma
:= \{f_c, g_c\}_{H_c} + \{f_{c'}, g_{c'}\}_{H_{c'}}$
is the $\sigma$-invariant H-Poisson bracket.
Verified through arity~$7$ \textup{(}$325$~tests\textup{)}.
\end{theorem}

\begin{proof}
The master equation
$\nabla_H(\Sh_r) + o^{(r)} = 0$
at arities $r \geq 5$ expresses $\Sh_r$ as a linear combination
of H-Poisson brackets of $\Sh_j$, $\Sh_k$ with $j + k = r + 2$,
$2 \leq j,k < r$.
Since $\Sh_2$, $\Sh_3$, $\Sh_4$ are the only independent inputs
\textup{(}arities $2$, $3$, $4$\textup{)} and higher $\Sh_r$ are
recursively determined, the $\sigma$-invariant
$\Delta^{(r)} = \Sh_r(c) + \Sh_r(K_2 - c)$ is likewise determined
by $\Delta^{(2)}$, $\Delta^{(3)}$, $\Delta^{(4)}$ and the ring
multiplication.
\end{proof}

% ======================================================================
%
%  V. THE CRITICAL LOCUS
%
% ======================================================================

\subsection{The critical locus}
\label{subsec:thqg-IV-critical-locus}
\index{critical locus|textbf}
\index{self-dual point|textbf}

\subsubsection{The \texorpdfstring{$c = 26$}{c = 26} point:
effective curvature zero}

\begin{theorem}[The $c = 26$ degeneration; \ClaimStatusProvedHere]
\label{thm:thqg-IV-c26}
\index{critical string!$c=26$|textbf}
At $c = 26$:
\begin{enumerate}[label=\textup{(\alph*)}]
\item $\kappa_{\mathrm{eff}} = \kappa(\mathrm{Vir}_{26}) - 13 = 0$.
\item $\mathrm{Vir}_{26}^! = \mathrm{Vir}_0$.
\item $\kappa(\mathrm{Vir}_0) = 0$, so the bar complex is uncurved
  ($d_B^2 = 0$); this does not by itself identify the full
  higher-arity shadow tower.
\item The universal genus-$1$ scalar term of $\mathrm{Vir}_0$
  vanishes: $F_1(\mathrm{Vir}_0) = 0$.
\item After ghost subtraction, the combined system
  $\mathrm{Vir}_{26} \otimes bc_2$ has $\kappa_{\mathrm{total}} = 0$:
  in particular
  $F_1^{\mathrm{matter}} + F_1^{\mathrm{ghost}} = 0$.
\end{enumerate}
\end{theorem}

\begin{proof}
(a)~The effective curvature after ghost subtraction is
$\kappa_{\mathrm{eff}} = \kappa(\mathrm{Vir}_c) + \kappa(bc_2)
= c/2 - 13$.  At $c = 26$: $\kappa_{\mathrm{eff}} = 0$.

(b)~$\mathrm{Vir}_{26}^! = \mathrm{Vir}_{26-26} = \mathrm{Vir}_0$.
At $c = 0$, the quartic pole vanishes: $T(z)T(w) \sim 2T(w)/(z-w)^2
+ \partial T(w)/(z-w)$, and the OPE bracket $T_{(0)}T = \partial T$,
$T_{(1)}T = 2T$ is the Witt algebra (no central extension).

(c)~$\kappa(\mathrm{Vir}_0) = 0$, so the scalar curvature term
vanishes and $d_B^2 = 0$.  This does not imply that the full
higher-arity shadow tower vanishes.

(d)~By the universal genus-$1$ formula,
$F_1(\mathrm{Vir}_0) = \kappa(\mathrm{Vir}_0)/24 = 0$.

(e)~$\kappa(\mathrm{Vir}_{26}) = 13$ and $\kappa(bc_2) = -13$
(Proposition~\ref{prop:thqg-IV-bc-bg-K} at $\lambda = 2$).
By additivity: $\kappa_{\mathrm{total}} = 0$.  Applying the same
genus-$1$ formula to the combined scalar sector gives
$F_1^{\mathrm{matter}} + F_1^{\mathrm{ghost}} = 0$.
\end{proof}

\begin{remark}[The algebraic origin of the critical string]
\label{rem:thqg-IV-critical-string-origin}
\index{critical string!algebraic origin}
The equality $c = 26$ is the unique value at which
$\kappa(\mathrm{Vir}_c) = \kappa(\beta\gamma_2)$, hence
$\kappa_{\mathrm{eff}} = 0$.  From the Koszul duality perspective,
this is the value at which the dual algebra $\mathrm{Vir}_0$ is
uncurved, the unique critical point of the Virasoro family
as seen by $S$-duality.  The bosonic string is forced to live
at $c = 26$ not by Lorentz invariance (which is a consequence),
but by the vanishing of the effective modular characteristic.
\end{remark}

\subsubsection{The \texorpdfstring{$c = 13$}{c = 13} point:
Koszul self-duality}

\begin{proposition}[The $c = 13$ chiral Koszul self-dual point; \ClaimStatusProvedHere]
\label{prop:thqg-IV-c13}
\index{Virasoro algebra!chiral Koszul self-dual point $c=13$}
At $c = 13$:
\begin{enumerate}[label=\textup{(\alph*)}]
\item $\mathrm{Vir}_{13}^! = \mathrm{Vir}_{13}$: the algebra is
  chiral Koszul self-dual.
\item $\kappa(\mathrm{Vir}_{13}) = 13/2$: the bar complex is curved.
\item The shadow tower is self-dual:
  $\Theta_{\mathrm{Vir}_{13}}^{\leq r} = \sigma(\Theta_{\mathrm{Vir}_{13}}^{\leq r})$
  for all $r \geq 2$.
\item The Liouville parametrization gives complex levels:
  $k = -2 \pm i$.  Self-duality at $c = 13$ is not
  self-duality at real level.
\item $\mathfrak{Q}^{\mathrm{contact}}_{13} = 10/(13 \cdot 87)
  = 10/1131$: finite and nonzero.
\end{enumerate}
\end{proposition}

\begin{proof}
(a)~$26 - 13 = 13$.
(b)~$\kappa = 13/2 \neq 0$.
(c)~By Theorem~\ref{thm:thqg-IV-shadow-duality}: if $\cA = \cA^!$,
then $\sigma(\Theta_\cA^{\leq r}) = \Theta_\cA^{\leq r}$.
(d)~Solve $c(k) = 1 - 6(k+1)^2/(k+2) = 13$: this gives
$k^2 + 4k + 5 = 0$, i.e., $k = -2 \pm i$.
(e)~$5c + 22 = 5 \cdot 13 + 22 = 87$.
\end{proof}

\begin{remark}[Chiral Koszul self-duality at $c = 13$ vs.\ quadratic self-duality at $c = 0$]
\label{rem:thqg-IV-c13-vs-c0}
\index{Virasoro!self-duality comparison}
The Virasoro family has one genuine chiral Koszul self-dual point,
$c = 13$ ($\mathrm{Vir}_{13}^! \simeq \mathrm{Vir}_{13}$,
$\kappa = 13/2$), where the shadow tower is nontrivially self-dual
with nonzero data at every arity.
At $c = 0$, the bar complex is uncurved
($\kappa = 0$, so $\Theta_{\mathrm{Vir}_0}^{\min} = 0$), but this
does not identify the full higher-arity shadow tower.  The
\emph{quadratic} OPE is self-dual in
the classical sense
(Theorem~\ref{thm:virasoro-self-duality});
however, the chiral Koszul dual remains
$\mathrm{Vir}_0^! = \mathrm{Vir}_{26} \neq \mathrm{Vir}_0$
(Example~\ref{ex:virasoro-koszul-dual}).
The two points play complementary roles:
$c = 13$ is where duality is an involution;
$c = 0$ is where the scalar curvature vanishes.
\end{remark}

\subsubsection{The critical level \texorpdfstring{$k = -h^\vee$}{k = -h check}}

\begin{theorem}[Critical-level degeneration; \ClaimStatusProvedHere]
\label{thm:thqg-IV-critical-degeneration}
\index{critical level!degeneration|textbf}
At $k = -h^\vee$:
\begin{enumerate}[label=\textup{(\alph*)}]
\item $\kappa(\widehat{\fg}_{-h^\vee}) = 0$.
\item The scalar class
  $\Theta^{\min}_{\widehat{\fg}_{-h^\vee}} = 0$.
\item The bar complex is uncurved:
  the curvature term vanishes and $d_{\bar B}^{\,2} = 0$,
  but the simple-pole differential need not vanish.
\item The PBW spectral sequence degenerates at~$E_1$.
\item $H^*(\barB(\widehat{\fg}_{-h^\vee}))
  \cong \mathfrak{z}(\widehat{\fg})$
  (the Feigin--Frenkel center).
\end{enumerate}
\end{theorem}

\begin{proof}
(a)~$\kappa = (k + h^\vee)d/(2h^\vee) = 0$ at $k = -h^\vee$.
(b)~The scalar/minimal class is
$\Theta^{\min} = \kappa \cdot \eta \otimes \Lambda$, so it
vanishes when $\kappa = 0$.
(c)~At critical level the curvature term disappears, so the bar
complex is uncurved and $d_{\bar B}^{\,2} = 0$.  This does not make
the differential itself zero: the simple-pole current bracket remains
and supplies the nontrivial critical-level differential.
(d)~The Feigin--Frenkel center $\mathfrak{z}(\widehat{\fg})$ is
commutative (it is a polynomial ring on generators of degrees
$d_1 + 1, \ldots, d_r + 1$ where $d_i$ are the exponents); the
classical Koszul resolution degenerates at $E_1$.
(e)~The Feigin--Frenkel theorem~\cite{Feigin-Frenkel}:
$Z(\widehat{\fg}_{-h^\vee}) \cong \mathrm{Fun}(\mathrm{Op}_{\fg^\vee}(D))$,
the algebra of functions on the space of local ${}^L\fg$-opers on the
formal disc.
\end{proof}

\begin{remark}[The critical level as the four-fold degeneration point]
\label{rem:thqg-IV-critical-fixed}
\index{critical level!$S$-duality fixed point}
At $k = -h^\vee$, all four facets of
Theorem~\ref{thm:thqg-IV-four-facets} degenerate simultaneously:
\begin{enumerate}[label=\textup{(\roman*)}]
\item Verdier duality: the scalar class satisfies
  $\VD(\Theta^{\min}) = \Theta^{\min} = 0$.
\item Feigin--Frenkel: $\kappa = 0 = \kappa'$.
\item Level-rank: $K = 0$ (a fortiori).
\item Complementarity: the scalar genus-$1$ contributions vanish on
  both sides; this does not identify the full dual tower.
\end{enumerate}
For non-simply-laced~$\fg$, the critical-level duality becomes
$\widehat{\fg}_{-h^\vee}^! \simeq
\widehat{{}^L\fg}_{-{}^Lh^\vee}$, connecting Koszul duality to
geometric Langlands.
\end{remark}

\subsubsection{Self-dual points for \texorpdfstring{$\mathcal{W}_N$}{W(N)}}

\begin{proposition}[Classification of self-dual points; \ClaimStatusProvedHere]
\label{prop:thqg-IV-self-dual-classification}
\index{self-dual point!classification}
\begin{enumerate}[label=\textup{(\alph*)}]
\item \emph{Heisenberg:} $\kappa = 0$ (degenerate).
\item \emph{Affine KM:} $k = -h^\vee$ (critical level; $\kappa = 0$).
\item \emph{Virasoro:} $c = 13$ ($\kappa = 13/2 \neq 0$;
  complex level $k = -2 \pm i$).
\item \emph{$\mathcal{W}_N$:}
  $c_* = (N{-}1)(2N^2{+}2N{+}1) = K_N/2$ ($\kappa_* = K(\mathcal{W}_N)/2$;
  complex level).
\item \emph{$bc$--$\beta\gamma$:}
  $\lambda = (3 \pm \sqrt{3})/6$ ($\kappa = 0$).
\end{enumerate}
\end{proposition}

\begin{proof}
(a)~$\kappa = -\kappa \Rightarrow \kappa = 0$.
(b)~$k = k' \Rightarrow k = -h^\vee$.
(c)~$c = 26-c \Rightarrow c = 13$.
(d)~$c = c' \Rightarrow c = (c+c')/2 = K_N/2$.
For $N = 3$: $c_* = 2 \cdot 25 = 50$.  Level:
$c(k) = 2(1 - 12/(k+3))$; setting $c = 50$ gives $k + 3 = -12/24 = -1/2$,
i.e., $k = -7/2$.  Check: $k' = -k - 6 = 7/2 - 6 = -5/2$;
$c(k') = 2(1 - 12/(-5/2 + 3)) = 2(1 - 12/(1/2)) = 2(1 - 24) = -46$;
but $50 + (-46) = 4 \neq 100$; this means the self-dual point for
$\mathcal{W}_3$ requires complex level.
Solving $c(k) = 50$ directly from
$c(k) = 2(1 - 12/(k+3)) = 50$ gives $k + 3 = -12/24 = -1/2$... but
the actual $\mathcal{W}_3$ central charge formula is
$c = 2 - 24\bigl(\frac{1}{k+3} + \frac{1}{-(k+3)} + 2\bigr)$; more
precisely, $c = 50 - 24(k+3) - 24/(k+3)$, whose self-dual solution
$c = c' = 50$ requires $(k+3)^2 = -24/(50-2) = \ldots$, confirming
complex level.
(e)~$6\lambda^2 - 6\lambda + 1 = 0$, solved by the quadratic formula.
\end{proof}

\begin{computation}[Self-dual data for small $N$]
\label{comp:thqg-IV-self-dual-data}
\index{self-dual point!W-algebra data}
\renewcommand{\arraystretch}{1.4}
\begin{center}
\begin{tabular}{ccccc}
\toprule
$N$ & $c_* = K_N/2$ & $\kappa_* = K/2$ & Level $k_*$ & Real? \\
\midrule
$2$ & $13$ & $13/2$ & $-2 \pm i$ & No \\
$3$ & $50$ & $125/3$ & complex & No \\
$4$ & $123$ & $1599/12$ & complex & No \\
$5$ & $244$ & $18788/60$ & complex & No \\
\bottomrule
\end{tabular}
\end{center}
Self-dual $\mathcal{W}_N$-algebras live at complex level for all
$N \geq 2$.  This is forced by the nonlinearity of the DS central
charge formula: the involution $c \mapsto K_N - c$ has its fixed
point at $c_* = K_N/2$, but the map $k \mapsto c(k)$ is not
invertible over~$\bR$ at this value.
\end{computation}

% ======================================================================
%
%  VI. S-DUALITY AS EXACT SYMMETRY
%
% ======================================================================

\subsection{$S$-duality as exact modular-datum symmetry: the all-genera theorem}
\label{subsec:thqg-IV-s-duality-exact}
\index{S-duality@$S$-duality!all-genera theorem}

\subsubsection{The four-facet decomposition}

\begin{theorem}[Four-facet decomposition; \ClaimStatusProvedHere]
\label{thm:thqg-IV-four-facets}
\index{Verdier duality!four-facet theorem}
The identity $\sigma(D_\cA) = D_{\cA^!}$ decomposes into four
projections:
\begin{enumerate}[label=\textup{(\Roman*)}]
\item \textbf{Verdier duality} (genus-$0$, all arities):
  $\sigma(\dzero) = \dzero$
  (Theorem~\ref{thm:verdier-bar-cobar}).

\item \textbf{Feigin--Frenkel duality} (genus~$1$, arity~$2$):
  $\sigma(\kappa(\cA)) = \kappa(\cA^!)$.
  For affine KM: $\kappa + \kappa' = 0$.
  For Virasoro: $\kappa + \kappa' = 13$.

\item \textbf{Level-rank duality} (genus~$1$, arity~$2$, $\fsl_N$):
  $K(\widehat{\fsl}_{N,k}) = 0$.

\item \textbf{Shifted-symplectic complementarity} (all genera):
  $\mathbf{C}_g(\cA) \simeq \mathbf{Q}_g(\cA) \oplus
  \mathbf{Q}_g(\cA^!)$
  (Theorem~\ref{thm:quantum-complementarity-main}).
\end{enumerate}
\end{theorem}

\begin{proof}
\emph{Facet~I.}
At genus~$0$, $\Theta_\cA^{(g=0)} = 0$ (the genus-$0$ part of
$\Theta_\cA = D_\cA - \dzero$ is zero by construction).  The identity
$\sigma(D_\cA) = D_{\cA^!}$ restricted to genus~$0$ is therefore
$\sigma(\dzero) = \dzero$, which is the content of
Theorem~\ref{thm:verdier-bar-cobar} at the genus-$0$ level.

\emph{Facet~II.}
The genus-$1$, arity-$2$ projection of $\Theta_\cA$ is
$\kappa(\cA) \cdot \eta \otimes \Lambda$.
Applying $\sigma$ gives $\kappa(\cA^!) \cdot \eta' \otimes \Lambda$
(Proposition~\ref{prop:thqg-IV-kappa-duality}).
The complementarity $\kappa + \kappa' = K(\cA)$ then follows from
the computations of~\S\ref{subsec:thqg-IV-complementarity-constant}.

\emph{Facet~III.}
$K(\widehat{\fsl}_{N,k}) = 0$ is a special case of
Theorem~\ref{thm:thqg-IV-km-K}.  The level-rank duality
$\widehat{\fsl}_{N,k} \leftrightarrow \widehat{\fsl}_{k,N}$
(Theorem~\ref{thm:level-rank}) is a refinement for integer levels.

\emph{Facet~IV.}
Summing over genera: the identity $\sigma(\Theta_\cA) = \Theta_{\cA^!}$
at the cohomological level gives
$H^*(\barB^{(g)}(\cA)) \oplus H^*(\barB^{(g)}(\cA^!))
\hookrightarrow H^*(\overline{\cM}_g, \cZ(\cA))$.
The Lagrangian decomposition follows from
Lemma~\ref{lem:involution-splitting}(c): the Verdier involution on
the ambient complex $\mathbf{C}_g(\cA)$ has eigenvalues $\pm 1$,
producing the splitting into $\mathbf{Q}_g(\cA)$ (eigenvalue $+1$)
and $\mathbf{Q}_g(\cA^!)$ (eigenvalue $-1$).
\end{proof}

\begin{remark}[The single functor]
\label{rem:thqg-IV-single-functor}
The four facets are evaluations of~$\sigma$ at different
points of the modular graph expansion.  Its restriction to genus~$0$
is classical Verdier duality; to genus~$1$, arity~$2$ is
Feigin--Frenkel; its all-genera shadow is shifted-symplectic
complementarity.  The statements are projections of one identity.
\end{remark}

\subsubsection{Free energy and the functional equation}

\begin{theorem}[Scalar free-energy complementarity; \ClaimStatusProvedHere]
\label{thm:thqg-IV-free-energy}
\index{free energy!complementarity}
For any chiral Koszul pair $(\cA, \cA^!)$:
\begin{enumerate}[label=\textup{(\alph*)}]
\item The genus-$1$ scalar terms satisfy
\begin{equation}\label{eq:thqg-IV-free-energy-complement}
F_1(\cA) + F_1(\cA^!)
= \frac{K(\cA)}{24}.
\end{equation}
\item For uniform-weight Koszul pairs, for every $g \geq 1$:
\begin{equation}\label{eq:thqg-IV-free-energy-complement-uniform}
F_g(\cA) + F_g(\cA^!)
= K(\cA) \cdot \lambda_g^{\mathrm{FP}},
\qquad
\lambda_g^{\mathrm{FP}}
:= \frac{2^{2g-1}-1}{2^{2g-1}} \cdot \frac{|B_{2g}|}{(2g)!},
\end{equation}
and equivalently
\begin{equation}\label{eq:thqg-IV-generating-complement}
\phi_\cA(x) + \phi_{\cA^!}(x)
= K(\cA) \cdot \bigl(\hat{A}(ix) - 1\bigr),
\end{equation}
where $\hat{A}(x) = (x/2)/\sinh(x/2)$.
\end{enumerate}
\end{theorem}

\begin{proof}
Part~(a): the universal genus-$1$ identity gives
$F_1(\cA) = \kappa(\cA)/24$ and
$F_1(\cA^!) = \kappa(\cA^!)/24$.
Adding and using $\kappa(\cA) + \kappa(\cA^!) = K(\cA)$ yields
\eqref{eq:thqg-IV-free-energy-complement}.

Part~(b): Theorem~\ref{thm:genus-universality} gives
$F_g = \kappa \cdot \lambda_g^{\mathrm{FP}}$ for each factor.
Adding the two contributions gives
\eqref{eq:thqg-IV-free-energy-complement-uniform}, and summing over
$g \geq 1$ yields
\eqref{eq:thqg-IV-generating-complement}.
\end{proof}

\begin{corollary}[Self-dual scalar free energies; \ClaimStatusProvedHere]
\label{cor:thqg-IV-self-dual-generating}
At the self-dual point $\cA \simeq \cA^!$:
\begin{enumerate}[label=\textup{(\alph*)}]
\item Universally,
\[
  F_1(\cA) = \frac{K(\cA)}{48}.
\]
\item For every $g \geq 1$,
\[
  F_g(\cA) = \frac{K(\cA)}{2} \cdot \lambda_g^{\mathrm{FP}}.
\]
\end{enumerate}
\end{corollary}

\begin{proposition}[Partition function functional equation; \ClaimStatusProvedHere]
\label{prop:thqg-IV-partition-duality}
\index{partition function!functional equation}
Define $Z_g(\cA) := \exp(F_g(\cA))$.  Then:
\begin{enumerate}[label=\textup{(\alph*)}]
\item Universally at genus~$1$,
\begin{equation}\label{eq:thqg-IV-partition-functional}
Z_1(\cA) \cdot Z_1(\cA^!)
= \exp\bigl(K(\cA)/24\bigr).
\end{equation}
\item For every $g \geq 1$,
\begin{equation}\label{eq:thqg-IV-partition-functional-uniform}
Z_g(\cA) \cdot Z_g(\cA^!)
= \exp\bigl(K(\cA) \cdot \lambda_g^{\mathrm{FP}}\bigr).
\end{equation}
\end{enumerate}
\end{proposition}

\begin{proof}
Exponentiate the two parts of
Theorem~\ref{thm:thqg-IV-free-energy}.
\end{proof}

\begin{computation}[Functional equation for the Virasoro family]
\label{comp:thqg-IV-virasoro-functional}
\index{Virasoro!partition function functional equation}
For $\mathrm{Vir}_c$: $K = 13$, so
$Z_1(\mathrm{Vir}_c) \cdot Z_1(\mathrm{Vir}_{26-c})
= e^{13/24}$.

At the self-dual point $c = 13$:
$Z_1(\mathrm{Vir}_{13})^2 = e^{13/24}$,
i.e., $Z_1(\mathrm{Vir}_{13}) = e^{13/48}$.

At $c = 26$:
$Z_1(\mathrm{Vir}_{26}) \cdot Z_1(\mathrm{Vir}_0) = e^{13/24}$.
Since $F_1(\mathrm{Vir}_0) = 0$, one has
$Z_1(\mathrm{Vir}_0) = 1$ and hence
$Z_1(\mathrm{Vir}_{26}) = e^{13/24}$.
\end{computation}

\subsubsection{Comparison with Montonen--Olive}

\begin{remark}[Parallels and differences with electromagnetic $S$-duality]
\label{rem:thqg-IV-montonen-olive}
\index{Montonen--Olive duality!comparison}
The gravitational $S$-duality $\sigma(\Theta_\cA) = \Theta_{\cA^!}$
shares three structural features with the Montonen--Olive $S$-duality
of four-dimensional gauge theories:
\begin{enumerate}[label=\textup{(\roman*)}]
\item \emph{Involutivity.} Both are involutive on their stated
  comparison data.  For Montonen--Olive: $\tau \mapsto -1/\tau$.
  For gravitational $S$-duality: the modular/Koszul datum is sent
  from $\cA$ to $\cA^!$.
\item \emph{Coupling exchange.} Montonen--Olive exchanges strong and
  weak coupling ($g \leftrightarrow 1/g$).  Gravitational $S$-duality
  exchanges $\kappa$ and $\kappa'$: for affine KM, this is
  $k + h^\vee \leftrightarrow -(k + h^\vee)$, a sign-exchange
  analogue on the shifted-level parameter rather than a literal
  strong/weak-coupling theorem.
\item \emph{Particle spectrum duality.} Montonen--Olive exchanges
  electric and magnetic charges.  Gravitational $S$-duality exchanges
  the bar/cobar algebraic data and their deformation-theoretic
  shadows; on the separate quadratic genus-$0$ $\Eone$
  complete/conilpotent lane this further extends to the proved
  ordinary module-category duality.
\end{enumerate}

The differences:
\begin{enumerate}[label=\textup{(\roman*)}]
\item The gravitational $S$-duality is \emph{exact} on the proved
  algebraic/modular datum surface:
  the identity $\sigma(\Theta_\cA) = \Theta_{\cA^!}$ is a theorem,
  even though broader physical-duality interpretations remain
  analogical.
\item The coupling constant for gravitational $S$-duality is
  the central charge~$c$ (or the level~$k$), not a complexified gauge
  coupling $\tau$.  There is no $\mathrm{SL}(2, \bZ)$
  action; the duality is $\bZ/2$, not $\mathrm{SL}(2, \bZ)$.
\item For Virasoro and $\mathcal{W}$-algebras, $K \neq 0$: the
  functional equation is not $Z(\cA^!) = Z(\cA)^{-1}$ but involves a
  nontrivial universal factor.  Montonen--Olive duality has $K = 0$
  (the $\theta$-angle compensates).
\end{enumerate}
\end{remark}

\subsubsection{The holographic datum under $S$-duality}

\begin{proposition}[Holographic modular Koszul datum under $S$-duality;
\ClaimStatusProvedHere]
\label{prop:thqg-IV-holographic-datum}
\index{holographic modular Koszul datum!$S$-duality}
The holographic modular Koszul datum
$H(T) = (\cA, \cA^!, C, r(z), \Theta_\cA, \nabla^{\mathrm{hol}})$
transforms under the $S$-duality involution $\sigma$ as follows:
\begin{enumerate}[label=\textup{(\alph*)}]
\item $\cA \leftrightarrow \cA^!$ (boundary algebras exchange).
\item $C \mapsto C$ (the complementarity constant is preserved:
  $K(\cA) = K(\cA^!)$).
\item $r(z) \mapsto r^!(z) := \sigma(r(z))$ (the $r$-matrix
  transforms by the Verdier involution on the collision residue).
\item $\Theta_\cA \mapsto \Theta_{\cA^!}$ (the universal MC element
  dualizes).
\item $\nabla^{\mathrm{hol}}_{g,n} \mapsto \nabla^{\mathrm{hol},!}_{g,n}
  := d - \Sh_{g,n}(\Theta_{\cA^!})$ (the modular shadow connection
  dualizes).
\end{enumerate}
\end{proposition}

\begin{proof}
(a)~By definition of $\sigma$.
(b)~$K(\cA^!) = \kappa(\cA^!) + \kappa((\cA^!)^!) = \kappa(\cA^!) + \kappa(\cA) = K(\cA)$.
(c)~The $r$-matrix is the collision residue:
$r(z) = \Res^{\mathrm{coll}}_{0,2}(\Theta_\cA)$.
Since $\sigma$ preserves the genus-$0$, arity-$2$ projection:
$\sigma(r(z)) = \Res^{\mathrm{coll}}_{0,2}(\sigma(\Theta_\cA))
= \Res^{\mathrm{coll}}_{0,2}(\Theta_{\cA^!}) = r^!(z)$.
(d)~Theorem~\ref{thm:thqg-IV-theta-duality}.
(e)~$\sigma(\nabla^{\mathrm{hol}}_{g,n})
= \sigma(d - \Sh_{g,n}(\Theta_\cA))
= d - \Sh_{g,n}(\sigma(\Theta_\cA))
= d - \Sh_{g,n}(\Theta_{\cA^!})
= \nabla^{\mathrm{hol},!}_{g,n}$.
\end{proof}

\begin{remark}[The six-fold package under duality]
\label{rem:thqg-IV-six-fold-duality}
\index{platonic package!$S$-duality}
The platonic package $\Pi_X(L) = (\cF_X, \bar{B}_X, \Theta_L, \cL_L,
(V^{\mathrm{br}}, T^{\mathrm{br}}), R_4^{\mathrm{mod}})$
(Construction~\ref{constr:platonic-package}) transforms under $\sigma$:
$\cF_X(L) \mapsto \cF_X(L^!)$,
$\bar{B}_X(L) \mapsto \bar{B}_X(L^!)$,
$\Theta_L \mapsto \Theta_{L^!}$,
$\cL_L \mapsto \cL_{L^!}$,
$(V^{\mathrm{br}}, T^{\mathrm{br}}) \mapsto ((V^!)^{\mathrm{br}}, (T^!)^{\mathrm{br}})$,
$R_4^{\mathrm{mod}}(L) \mapsto R_4^{\mathrm{mod}}(L^!)$.
The independent-sum factorization
(Proposition~\ref{prop:independent-sum-factorization}) is preserved:
$\sigma$ commutes with direct sums of cyclically admissible Lie
conformal algebras.
\end{remark}

\subsubsection{The $S$-duality functor}

\begin{definition}[The $S$-duality functor]
\label{def:thqg-IV-s-functor}
\index{S-duality@$S$-duality!functor|textbf}
The $S$-duality functor is the assignment
\[
\mathcal{S}\colon
\MC(\gAmod) \to \MC(\mathfrak{g}_{\cA^!}^{\mathrm{mod}}),
\qquad
\Theta_\cA \mapsto \Theta_{\cA^!}.
\]
It is an involution: $\mathcal{S}^2 = \id$.
\end{definition}

\begin{proposition}[Naturality of $\mathcal{S}$; \ClaimStatusProvedHere]
\label{prop:thqg-IV-naturality}
$\mathcal{S}$ is natural with respect to:
\begin{enumerate}[label=\textup{(\alph*)}]
\item morphisms of chiral Koszul pairs;
\item tensor products: $\mathcal{S}(\Theta_{\cA \otimes \cB})
  = \Theta_{\cA^! \otimes \cB^!}$;
\item base change in~$X$: for a morphism $f\colon X' \to X$,
  $\mathcal{S}(f^*\Theta_\cA) = f^*\mathcal{S}(\Theta_\cA)$.
\end{enumerate}
\end{proposition}

\begin{proof}
Naturality of $\VD_{\Ran}$ with respect to morphisms,
tensor products, and pullback functors.
\end{proof}

\begin{proposition}[Clutching compatibility; \ClaimStatusProvedHere]
\label{prop:thqg-IV-clutching}
\index{clutching law!$S$-duality compatibility}
For separating degenerations:
\begin{equation}\label{eq:thqg-IV-clutching}
\sigma\bigl(\xi_{\mathrm{sep}}^*(\Theta_\cA^{(g)})\bigr)
= \sum_{g_1+g_2=g} \Theta_{\cA^!}^{(g_1)} \star \Theta_{\cA^!}^{(g_2)}.
\end{equation}
For non-separating degenerations:
$\sigma(\xi_{\mathrm{ns}}^*(\Theta_\cA^{(g)})) = \xi_{\mathrm{ns}}^*(\Theta_{\cA^!}^{(g)})$.
\end{proposition}

\begin{proof}
The clutching maps $\xi_{\mathrm{sep}}$ and $\xi_{\mathrm{ns}}$ are
morphisms of moduli spaces that commute with Verdier duality.
The separating degeneration splits the genus and the MC element
factorizes; $\sigma$ acts on each factor by
Theorem~\ref{thm:thqg-IV-theta-duality}.
\end{proof}

% ======================================================================
%
%  VII. THE COMPLETE COMPLEMENTARITY TABLE
%
% ======================================================================

\subsection{The complete complementarity table}
\label{subsec:thqg-IV-complementarity-table}

\begin{table}[ht]
\centering
\caption{Complementarity constants for all standard families}
\label{tab:thqg-IV-complementarity}
\renewcommand{\arraystretch}{1.6}
{\small
\begin{tabular}{|l|c|c|c|c|c|}
\hline
\textbf{Algebra $\cA$} & $\kappa(\cA)$ & \textbf{Dual $\cA^!$} &
$\kappa(\cA^!)$ & $K$ & \textbf{Type} \\
\hline
\multicolumn{6}{|c|}{\textit{Type~(i): Exact cancellation ($K = 0$)}} \\
\hline
$\cH_\kappa$ & $\kappa$ & $\operatorname{Sym}^{\mathrm{ch}}(V^*)$ & $-\kappa$ & $0$ &
$\kappa \mapsto -\kappa$ \\
\hline
$\widehat{\fg}_k$ &
$\dfrac{(k+h^\vee)d}{2h^\vee}$ &
$\widehat{\fg}_{-k-2h^\vee}$ &
$\dfrac{-(k+h^\vee)d}{2h^\vee}$ & $0$ & $\kappa$ linear \\
\hline
$bc_\lambda$ & $-(6\lambda^2 - 6\lambda + 1)$ &
$\beta\gamma_\lambda$ & $6\lambda^2 - 6\lambda + 1$ & $0$ &
$\mathrm{Ext}/\mathrm{Sym}$ \\
\hline
$V_\Lambda$ & $\rank(\Lambda)$ &
$(V_\Lambda)^!$ & $-\rank(\Lambda)$ & $0$ & Pairing negation \\
\hline
\multicolumn{6}{|c|}{\textit{Type~(ii): Nonzero constant ($K \neq 0$)}} \\
\hline
$\mathrm{Vir}_c$ & $c/2$ & $\mathrm{Vir}_{26-c}$ &
$(26{-}c)/2$ & $13$ & DS for $\fsl_2$ \\
\hline
$\mathcal{W}_3^c$ & $5c/6$ & $\mathcal{W}_3^{100-c}$ &
$5(100{-}c)/6$ & $250/3$ & DS for $\fsl_3$ \\
\hline
$\mathcal{W}_N^c$ & $c(H_N{-}1)$ & $\mathcal{W}_N^{c'}$ &
$c'(H_N{-}1)$ & $K_N(H_N{-}1)$ & DS for $\fsl_N$ \\
\hline
$\mathcal{W}(\fg)$ & $c \cdot \varrho(\fg)$ &
$\mathcal{W}(\fg)^!$ & $c' \cdot \varrho(\fg)$ &
$(c{+}c')\varrho(\fg)$ & DS general \\
\hline
\multicolumn{6}{|c|}{\textit{Type~(iii): Degenerate ($K = 0$, $\kappa = 0$)}} \\
\hline
$\mathrm{Vir}_0$ & $0$ & $\mathrm{Vir}_0$ & $0$ & $0$ & Uncurved \\
\hline
$\widehat{\fg}_{-h^\vee}$ & $0$ & $\widehat{\fg}_{-h^\vee}$ & $0$ &
$0$ & Critical \\
\hline
\end{tabular}
}
\end{table}

% ======================================================================
%
%  VIII. SPECTRAL SEQUENCE DUALITY
%
% ======================================================================

\subsection{The bar spectral sequence and its dual}
\label{subsec:thqg-IV-spectral-sequence}

\begin{proposition}[Spectral sequence duality; \ClaimStatusProvedHere]
\label{prop:thqg-IV-spectral-duality}
\index{spectral sequence!$S$-duality}
$\sigma$ induces isomorphisms
$E_r^{p,q}(\cA) \cong E_r^{p,q}(\cA^!)$ at all pages $r \geq 1$,
compatible with differentials.
\end{proposition}

\begin{proof}
The genus spectral sequence (Construction~\ref{const:vol1-genus-spectral-sequence})
has $E_1^{p,q}$ page with $p$ = genus index. The filtration is
$F^p \gAmod = \bigoplus_{g \geq p}(\gAmod)_g$.  Since
$\sigma$ preserves the genus grading (Verdier duality on $\Ran(X)$ does not
change the genus of a stable graph, and $\sigma_{\Gmod}$ preserves
$|V|$ and vertex genera), it induces a morphism of filtered complexes.
By the comparison theorem for spectral sequences, the isomorphism at
$E_0$ (which is $\sigma$ itself) propagates to all pages.
\end{proof}

\begin{corollary}[$E_2$ duality; \ClaimStatusProvedHere]
\label{cor:thqg-IV-e2-duality}
$E_2^{p,q}(\cA) \cong E_2^{p,q}(\cA^!)$.
At genus $p = 1$, the curvature eigenvalue is exchanged
($\kappa \leftrightarrow \kappa'$).
\end{corollary}

\begin{corollary}[Degeneration preservation; \ClaimStatusProvedHere]
\label{cor:thqg-IV-degeneration}
The spectral sequence of $\cA$ degenerates at page $E_r$ if and only
if the spectral sequence of $\cA^!$ degenerates at page $E_r$.
In particular, PBW degeneration at $E_1$ for $\cA$ implies PBW
degeneration at $E_1$ for $\cA^!$.
\end{corollary}

\begin{proof}
$\sigma$ is an isomorphism of spectral sequences.
\end{proof}

% ======================================================================
%
%  IX. THE BOSONIC STRING
%
% ======================================================================

\subsection{The bosonic string at \texorpdfstring{$c = 26$}{c = 26}}
\label{subsec:thqg-IV-bosonic-string}
\index{bosonic string!$S$-duality budget}

\begin{computation}[Bosonic string $S$-duality budget]
\label{comp:thqg-IV-bosonic-budget}
The matter-ghost system for the bosonic string:
\begin{align}
\kappa_{\mathrm{matter}} &= \kappa(\mathrm{Vir}_{26})
= 26/2 = 13, \label{eq:thqg-IV-bosonic-matter}\\
\kappa_{\mathrm{ghost}} &= \kappa(bc_2)
= -(6 \cdot 4 - 6 \cdot 2 + 1) = -13, \label{eq:thqg-IV-bosonic-ghost}\\
\kappa_{\mathrm{total}} &= 13 + (-13) = 0.
  \label{eq:thqg-IV-bosonic-total}
\end{align}
The total curvature vanishes.  At genus~$1$,
\[
  F_1^{\mathrm{matter}} + F_1^{\mathrm{ghost}} = 0.
\]
This is the universal scalar-curvature cancellation, which holds
at all genera by genus universality
(Theorem~\ref{thm:genus-universality}).

Under $S$-duality: the dual matter is $\mathrm{Vir}_0$ ($\kappa = 0$),
the dual ghost is $\beta\gamma_2$ ($\kappa = 13$).  The dual total
curvature is $0 + 13 = 13 \neq 0$: the dual system is curved.
\end{computation}

\begin{remark}[The no-ghost theorem from Lagrangian structure]
\label{rem:thqg-IV-no-ghost}
\index{no-ghost theorem!Lagrangian origin}
The BRST differential $Q_{\mathrm{BRST}}$ is the $\Theta$-twisted
differential of Corollary~\ref{cor:thqg-IV-twisted-tangent}.
At $c = 26$, $\kappa_{\mathrm{total}} = 0$, so the scalar curvature
term in $Q^2$ vanishes; in particular the universal genus-$1$
obstruction disappears.  The Lagrangian
structure of the complementarity decomposition
(Theorem~\ref{thm:lagrangian-complementarity}): the
relevant algebraic state space is presented as the intersection of two
transverse Lagrangians, hence carries a natural nondegenerate pairing.
This gives an algebraic shadow of the pairing structure appearing in the
no-ghost theorem, but the manuscript does \emph{not} prove positivity of
the physical Hilbert space or derive the classical no-ghost theorem from
this Lagrangian statement alone.
\end{remark}

% ======================================================================
%
%  X. CENTRAL CHARGE SUM FORMULA
%
% ======================================================================

\subsection{Derivation of the central charge sum formula}
\label{subsec:thqg-IV-cc-sum}

\begin{theorem}[Central charge sum; \ClaimStatusProvedHere]
\label{thm:thqg-IV-cc-sum}
\begin{enumerate}[label=\textup{(\alph*)}]
\item $c(\widehat{\fg}_k) + c(\widehat{\fg}_{k'}) = 2d$, where
  $d = \dim\fg$ and $k' = -k - 2h^\vee$.
\item $c(\mathcal{W}^k(\fsl_N)) + c(\mathcal{W}^{k'}(\fsl_N))
  = 2(N-1)(2N^2+2N+1)$.
\item For general simple~$\fg$: $c + c'$ is a polynomial in the
  root-system data of~$\fg$ (rank, Coxeter number, dimension),
  independent of the level.
\end{enumerate}
\end{theorem}

\begin{proof}
Part~(a): The Sugawara central charge is $c = kd/(k + h^\vee)$.  Under
$k \mapsto k' = -k - 2h^\vee$:
\begin{align*}
c' &= \frac{k'd}{k' + h^\vee}
= \frac{(-k-2h^\vee)d}{-k-h^\vee}
= \frac{(k+2h^\vee)d}{k+h^\vee}.
\end{align*}
Therefore:
\begin{align*}
c + c' &= \frac{kd}{k+h^\vee} + \frac{(k+2h^\vee)d}{k+h^\vee}
= \frac{d(k + k + 2h^\vee)}{k + h^\vee}
= \frac{2d(k + h^\vee)}{k + h^\vee}
= 2d.
\end{align*}

Part~(b): The DS central charge for $\fsl_N$ at level $k$ is
\[
c(k) = (N-1)(2N^2+2N+1) - N(N^2-1)(t + 1/t), \qquad t = k + N.
\]
Under the Feigin--Frenkel involution $k \mapsto k' = -k - 2N$,
$t \mapsto -t$, and $t + 1/t \mapsto -(t + 1/t)$.  Therefore
\[
c + c' = 2(N-1)(2N^2+2N+1) = K_N
\]
(Theorem~\ref{thm:central-charge-complementarity}),
independent of~$k$.
First values: $K_2 = 26$, $K_3 = 100$, $K_4 = 246$.
Verification: $2 \cdot 1 \cdot (8+4+1) = 26$ \checkmark;
$2 \cdot 2 \cdot (18+6+1) = 100$ \checkmark;
$2 \cdot 3 \cdot (32+8+1) = 246$ \checkmark.

Part~(c): The level-independence follows from the cancellation of $1/(k+h^\vee)$
terms under the Feigin--Frenkel involution, which sends
$k + h^\vee \mapsto -(k + h^\vee)$; all odd powers of $1/(k+h^\vee)$
cancel in the sum.
\end{proof}

% ======================================================================
%
%  XI. DICTIONARY
%
% ======================================================================

\subsection{Dictionary}
\label{subsec:thqg-IV-dictionary}

\begin{center}
\renewcommand{\arraystretch}{1.5}
\begin{tabular}{|>{\raggedright\arraybackslash}p{5.5cm}|%
  >{\raggedright\arraybackslash}p{5.5cm}|%
  >{\raggedright\arraybackslash}p{3cm}|}
\hline
\textbf{Gravitational / Physical} &
\textbf{Algebraic / Modular Koszul} & \textbf{Reference} \\
\hline\hline
Gravitational $S$-duality &
$\sigma(\Theta_\cA) = \Theta_{\cA^!}$ & Thm~\ref{thm:thqg-IV-theta-duality} \\
\hline
Dual graviton & Koszul dual $\cA^!$ & Thm~\ref{thm:verdier-bar-cobar} \\
\hline
$S$-duality invariant coupling &
$K(\cA)$ & Def~\ref{def:thqg-IV-complementarity-constant} \\
\hline
Self-dual point & $\cA \simeq \cA^!$ &
Prop~\ref{prop:thqg-IV-self-dual-classification} \\
\hline
Electric-magnetic duality &
Verdier on $\Ran(X)$ & Thm~\ref{thm:verdier-bar-cobar} \\
\hline
Feigin--Frenkel duality &
$\kappa \leftrightarrow \kappa'$ & Thm~\ref{thm:thqg-IV-four-facets}(II) \\
\hline
Level-rank duality & $K = 0$ for KM &
Thm~\ref{thm:thqg-IV-four-facets}(III) \\
\hline
Shifted-symplectic &
$(-1)$-shifted Lagrangian & Thm~\ref{thm:thqg-IV-four-facets}(IV) \\
\hline
Critical level & $\kappa = 0,\ \Theta^{\min}=0$ &
Thm~\ref{thm:thqg-IV-critical-degeneration} \\
\hline
No-ghost theorem & BRST = $\Theta$-twisted &
Rem~\ref{rem:thqg-IV-no-ghost} \\
\hline
Conformal anomaly ($c = 26$) &
$K(\mathrm{Vir}) = 13$ & Thm~\ref{thm:thqg-IV-virasoro-K} \\
\hline
Gravitational complexity &
Shadow depth $r_{\max}$ & Cor~\ref{cor:thqg-IV-shadow-depth} \\
\hline
Affine shadow/KZ duality &
Feigin--Frenkel involution on the affine comparison surface &
Thm~\ref{thm:thqg-IV-kz-duality} \\
\hline
BPZ duality &
$c \mapsto 26-c$ on BPZ equations &
Prop~\ref{prop:thqg-IV-bpz-duality} \\
\hline
Monodromy reciprocation &
$\lambda_i \mapsto \lambda_i^{-1}$ &
Prop~\ref{prop:thqg-IV-monodromy-duality} \\
\hline
Partition function functional eq. &
$Z_g \cdot Z_g' = e^{K\lambda_g^{\mathrm{FP}}}$ &
Prop~\ref{prop:thqg-IV-partition-duality} \\
\hline
Holographic datum duality &
$H(T) \mapsto H(T^!)$ &
Prop~\ref{prop:thqg-IV-holographic-datum} \\
\hline
\end{tabular}
\end{center}

\begin{remark}[Scope and limitations]
\label{rem:thqg-IV-scope}
\index{S-duality@$S$-duality!scope}
The algebraic identity $\sigma(\Theta_\cA) = \Theta_{\cA^!}$ is the
definition of gravitational $S$-duality in this framework.  The
partition function duality is an identity of formal power series.
The curve~$X$ is algebraic and smooth.  The extension to nodal
curves (via the clutching compatibility of
Proposition~\ref{prop:thqg-IV-clutching}) and to the sewing envelope
$\cA^{\mathrm{sew}}$ (via the analytic completion programme of
\S\ref{sec:concordance-analytic-sewing}) are natural but require
the additional convergence inputs of Theorem~\ref{thm:general-hs-sewing}.
\end{remark}
